\documentclass[11pt]{article}

% Change "review" to "final" to generate the final (sometimes called camera-ready) version.
% Change to "preprint" to generate a non-anonymous version with page numbers.
\usepackage[review]{acl}

% Standard package includes
\usepackage{times}
\usepackage{latexsym}

% For proper rendering and hyphenation of words containing Latin characters (including in bib files)
\usepackage[T1]{fontenc}
% For Vietnamese characters
% \usepackage[T5]{fontenc}
% See https://www.latex-project.org/help/documentation/encguide.pdf for other character sets

% This assumes your files are encoded as UTF8
\usepackage[utf8]{inputenc}

% This is not strictly necessary, and may be commented out,
% but it will improve the layout of the manuscript,
% and will typically save some space.
\usepackage{microtype}

% This is also not strictly necessary, and may be commented out.
% However, it will improve the aesthetics of text in
% the typewriter font.
\usepackage{inconsolata}

%Including images in your LaTeX document requires adding
%additional package(s)
\usepackage{graphicx}

\usepackage{amsmath}      % For align, align*, and \text{}
\usepackage{amssymb}      % For mathematical symbols and \mathbb
\usepackage{amsfonts}     % For math fonts
\usepackage{bm}           % For bold math symbols if needed
\usepackage{algorithm}    % For the Algorithm environment
\usepackage{algorithmic}  % For the algorithmic steps (IF, FOR, etc.)
\usepackage{float}        % For [H] exact figure placement

% If the title and author information does not fit in the area allocated, uncomment the following
%
%\setlength\titlebox{<dim>}
%
% and set <dim> to something 5cm or larger.


%% MY PACKEGES
\usepackage{tabularx}
\usepackage[table]{xcolor}
\usepackage{multirow}

\usepackage{makecell} % Позволяет делать переносы строк в ячейках
\newcolumntype{Y}{>{\centering\arraybackslash}X}


%% для аппендикса, раздел примеры
\usepackage[utf8]{inputenc}
\usepackage[T1]{fontenc}
\usepackage{xcolor}
\usepackage[skins,breakable]{tcolorbox}
\definecolor{dustyblue}{HTML}{728FA5}
% Настраиваем красивую рамку для примеров
\newtcolorbox{genbox}[1][]{
  enhanced,
  breakable,
  colback=dustyblue!10, % Светлая пыльно-синяя заливка фона
  colframe=dustyblue!90!black, % Темная пыльно-синяя рамка
  boxrule=0.5mm,
  arc=3mm,
  left=3mm, right=3mm, top=2mm, bottom=2mm,
  fonttitle=\bfseries\sffamily,
  coltitle=white,
  title={#1},
  shadow={2mm}{-1mm}{0mm}{black!20}
}
\newcommand{\loopmark}{\textcolor{dustyblue!80!black}{\textit{[collapsed loop]}}}


% \title{AdaPop: Adaptive Popularity improves  Unlearning in LLMs}
\title{The More Popular, The Harder to Forget: Adaptive Popularity for LLM Unlearning}

% Author information can be set in various styles:
% For several authors from the same institution:
% \author{Author 1 \and ... \and Author n \\
%         Address line \\ ... \\ Address line}
% if the names do not fit well on one line use
%         Author 1 \\ {\bf Author 2} \\ ... \\ {\bf Author n} \\
% For authors from different institutions:
% \author{Author 1 \\ Address line \\  ... \\ Address line
%         \And  ... \And
%         Author n \\ Address line \\ ... \\ Address line}
% To start a separate ``row'' of authors use \AND, as in
% \author{Author 1 \\ Address line \\  ... \\ Address line
%         \AND
%         Author 2 \\ Address line \\ ... \\ Address line \And
%         Author 3 \\ Address line \\ ... \\ Address line}

\author{First Author \\
  Affiliation / Address line 1 \\
  Affiliation / Address line 2 \\
  Affiliation / Address line 3 \\
  \texttt{email@domain} \\\And
  Second Author \\
  Affiliation / Address line 1 \\
  Affiliation / Address line 2 \\
  Affiliation / Address line 3 \\
  \texttt{email@domain} \\}

%\author{
%  \textbf{First Author\textsuperscript{1}},
%  \textbf{Second Author\textsuperscript{1,2}},
%  \textbf{Third T. Author\textsuperscript{1}},
%  \textbf{Fourth Author\textsuperscript{1}},
%\\
%  \textbf{Fifth Author\textsuperscript{1,2}},
%  \textbf{Sixth Author\textsuperscript{1}},
%  \textbf{Seventh Author\textsuperscript{1}},
%  \textbf{Eighth Author \textsuperscript{1,2,3,4}},
%\\
%  \textbf{Ninth Author\textsuperscript{1}},
%  \textbf{Tenth Author\textsuperscript{1}},
%  \textbf{Eleventh E. Author\textsuperscript{1,2,3,4,5}},
%  \textbf{Twelfth Author\textsuperscript{1}},
%\\
%  \textbf{Thirteenth Author\textsuperscript{3}},
%  \textbf{Fourteenth F. Author\textsuperscript{2,4}},
%  \textbf{Fifteenth Author\textsuperscript{1}},
%  \textbf{Sixteenth Author\textsuperscript{1}},
%\\
%  \textbf{Seventeenth S. Author\textsuperscript{4,5}},
%  \textbf{Eighteenth Author\textsuperscript{3,4}},
%  \textbf{Nineteenth N. Author\textsuperscript{2,5}},
%  \textbf{Twentieth Author\textsuperscript{1}}
%\\
%\\
%  \textsuperscript{1}Affiliation 1,
%  \textsuperscript{2}Affiliation 2,
%  \textsuperscript{3}Affiliation 3,
%  \textsuperscript{4}Affiliation 4,
%  \textsuperscript{5}Affiliation 5
%\\
%  \small{
%    \textbf{Correspondence:} \href{mailto:email@domain}{email@domain}
%  }
%}

\begin{document}
\maketitle











% \begin{abstract}
% LLM unlearning algorithms treat all facts in the forget set as equally hard to erase, ignoring that popular facts are encountered far more often during pretraining and are could be far more resistant to erasure. Existing gradient-based methods weight updates by per-token model confidence, a local signal that conflates surface fluency with parametric memorisation depth. We present \textbf{AdaPop} (Adaptive Popularity), which scales the unlearning signal by a popularity-dependent exponent derived from external knowledge bases (e.g., Wikidata), %\texttt{pop\_sum},
% and automates the forget-retain balance via a dual-ascent controller. Evaluated against a battery of SOTA unlearning methods, such as Weighted Gradient Ascent (WGA), %GA, GD, NPO, and WGA
% on three model families and two real-world benchmarks, AdaPop reduces forget ROUGE-L below $0.051$ while retaining ROUGE-L above $0.855$, and limits popular-paraphrase forget ROUGE-L below $0.060$ (a $3.5\times$ improvement over WGA). Internal metrics confirm deeper parametric erasure ($\Delta$Rank $+53{,}760$ vs.\ WGA $-3{,}714$) with MMLU preserved within $0.05$ of the pre-unlearning checkpoint.
% \end{abstract}

\begin{abstract}
%LLM unlearning algorithms treat all facts in the forget set as equally hard to erase, ignoring that popular facts are encountered far more often during pretraining and are consequently far more resistant to removal. Existing gradient-based unlearning techniques weight updates by per-token model confidence, a local signal that conflates surface fluency with parametric memorisation depth. We present \textbf{AdaPop} (Adaptive Popularity), which scales the unlearning signal by a popularity-dependent exponent derived from external knowledge bases, and automates the forget-retain balance via a dual-ascent controller. Evaluated against SOTA unlearning algorithms on three LLM families and two benchmarks DUET, RWKU, \textbf{AdaPop} consistently outperforms all non-collapsing baselines on paraphrase-generalised erasure and maintains a reliable forget-retain tradeoff across all model families. %Competing methods are at best narrowly competitive on surface metrics for one of three models and do not generalise to paraphrase evaluation. 

LLM unlearning techniques typically treat every example in a forget set as equally easy or hard to remove, even though frequently encountered (popular) facts are far more resistant to parametric removal. Existing gradient-based approaches weight token updates by model confidence, which is a local, surface-level measure that can conflate fluency with deep memorisation. We propose \textbf{AdaPop} (Adaptive Popularity): a robust approach that scales per-fact unlearning by a popularity-dependent exponent derived from an external proxy (e.g., Wikidata citelinks) and automatically balances forgetting vs. retention via a dual-ascent controller. Across three model families and two real-world knowledge benchmarks, \textbf{AdaPop} substantially improves paraphrase-generalised unlearning %and internal-representation disruption 
while preserving downstream capabilities. %Internal metrics and ablations show the popularity exponent and the dual controller are complementary components that together achieve deeper, more generalised erasure than competitive baselines.
Internal representation metrics confirm genuine parametric erasure rather than surface suppression, with general capability preserved across all models.
\end{abstract}



\section{Introduction}
\label{sec:intro}
% %% Describe what is the problem current algos (stability isses, popularity influence): Gradient Ascent, Gradient Difference, NPO, WGA;
%% Explain how we solve this popularity gap 
%% Which benchmark we use to evaluate (we use real-world facts QA benchmarks for solving real-words tasks; we scoring facts based on they popularity. if you have ascess to pretrain set, you can use information from it to scoring populaeiry; otherwise use 'approximate' metrics from large data corpus like wikidata to score facts. here popularity is the global metric for for unlerning regularization; the local (model) metric is log probes and its to factors influence to unlearning.) also for trade-off stability we automatically balance retain, that directly effect to forgetting.

The necessity of removing specific knowledge from Large Language Models (LLMs) has led to the development of various weight-updating unlearning algorithms. However, established methods often struggle with a fundamental stability trade-off between forgetting target information and preserving general model capabilities. Traditional Gradient Ascent (GA) simply maximizes the loss on the forget set, which frequently leads to catastrophic collapse of the model's linguistic coherence. Gradient Difference (GD) attempts to mitigate this by adding a retain loss, yet it still treats every token in the forget set as having equal importance, regardless of whether a token is a critical factual entity or a high-frequency noise word like an article or preposition. More recent approaches such as Negative Preference Optimization (NPO) and Weighted Gradient Ascent (WGA) attempt to refine the unlearning process by using the model's own internal states. NPO treats unlearning as a preference alignment problem where the model is pushed away from the "preferred" response. WGA uses the model's log-probabilities to weight the gradient updates, theoretically focusing the unlearning effort on tokens where the model is most confident. While these methods are more sophisticated than simple ascent, they rely exclusively on the model's current confidence as a local metric. Our core assumption is that log-probabilities alone do not always accurately reflect how deeply a fact is embedded in the model. We argue that incorporating an external proxy for "fact popularity" better conveys the necessary unlearning strength than internal model confidence alone.

To address this, we introduce AdaPop (Adaptive Popularity), an algorithm designed to make the unlearning process more flexible by modulating updates based on the global prevalence of a fact. The relationship between a fact's presence in a massive pretraining corpus and its resistance to unlearning is often non-linear. Popularity can be quantified directly if the pretraining data is accessible by measuring citelinks or co-occurrence frequencies. In cases where the original pretraining corpus is private or unknown, we utilize large-scale public data proxies. This includes structural knowledge bases like Wikidata or massive web-scale corpora like Common Crawl to estimate how often a subject-object relationship appears in the general distribution of human knowledge.

AdaPop uses a popularity-dependent exponent to adjust the influence of each fact during the optimization process. This ensures that the gradient updates are scaled appropriately: deeply entrenched, popular facts are handled with a different weighting logic than obscure, long-tail information that might only appear a handful of times in the training set. Furthermore, AdaPop implements a dual-ascent controller to automatically manage the balance between forgetting and retaining. This removes the need for manual hyperparameter tuning of the retain coefficient $\alpha$, allowing the algorithm to adaptively maintain model utility based on real-time drift statistics. We demonstrate the effectiveness of this popularity-aware approach on benchmarks like DUET and RWKU, where diverse fact distributions highlight the limitations of purely confidence-based methods.
%% VERSION 2
Machine unlearning~\cite{cao2015towards, bourtoule2021machine} aims to remove specific, targeted knowledge from a trained model without full retraining, enabling downstream control over privacy, safety, and compliance risks that arise when large language models (LLMs) memorize and re-produce training data~\cite{carlini2021extracting, tirumala2022memorization}. A common setup is targeted factual unlearning for QA or knowledge-grounded generation: given a ``forget set'' of (query, answer) pairs, the goal is to eliminate the model’s ability to produce those answers while preserving its general language and reasoning capabilities on held-out (``retain'') data. Existing solutions fall into a few families.
%Machine unlearning removes specific knowledge from a trained model without retraining from scratch~\cite{cao2015towards, bourtoule2021machine}, and has become practically necessary as LLMs are shown to memorize and regenerate personally identifiable information from pretraining data~\cite{carlini2021extracting, tirumala2022memorization}. 
Existing methods mainly apply gradient-based and preference-learning objectives to suppress target knowledge~\cite{jang2022knowledge, liu2022continual, zhang2024negative, wang2025rethinking}, and typically treat all forget-set data points as equally hard to erase. Facts encountered more frequently during pretraining are harder to remove than rare ones~\cite{kandpal2023large, merullo2025linear}; applying the same gradient force to all forget-set facts regardless produces either over-erasure of rare facts or under-erasure of popular ones. %This failure mode, the \emph{popularity gap}, leads to catastrophic or ineffective forgetting. Recent research papers~\cite{krishnannot, anna2026anatomy} empirically document it across multiple methods and benchmarks, but neither proposes a training method that addresses it.
Recent empirical work documents a persistent \emph{popularity gap}: facts encountered more frequently during pretraining are systematically harder to erase and remain extractable under adversarial prompts~\cite{krishnannot, anna2026anatomy}. These findings suggest that methods which rely purely on local, model-intrinsic signals are insufficient for robust, generalisable unlearning.

% \begin{figure}
%     \centering
%     \includegraphics[width=1\linewidth]{latex//img/sq_crop_adapop.png}
%     \caption{AdaPop assigns a per-fact exponent $\beta_i$ derived from an external popularity score, applying weaker unlearning to rare facts and stronger unlearning to popular ones. A dual-ascent controller adjusts the retain coefficient $\alpha$ at each epoch to keep retain-set loss within a preset tolerance.}
%     \label{fig:placeholder}
% \end{figure}

\begin{figure*}[t]
    \centering
    \includegraphics[width=0.85\textwidth]{latex/img/crop_adapop.pdf}
    %\caption{AdaPop assigns a per-fact exponent $\beta_i$ derived from an external popularity score, applying weaker unlearning to rare facts and stronger unlearning to popular ones. A dual-ascent controller adjusts the retain coefficient $\alpha$ at each epoch to keep retain-set loss within a preset tolerance.}
    \caption{Overview of AdaPop. The popularity exponent $\beta_i$ reweights per-fact token weights; the dual-ascent controller updates $\alpha$ each epoch to keep the retain loss within a preset tolerance $\varepsilon$.}
    \label{fig:adapop_workflow}
\end{figure*}

To address this issue, we introduce \textbf{AdaPop} (\textbf{Ada}ptive \textbf{Pop}ularity), a popularity-weighted adaptive gradient unlearning method (Figure~\ref{fig:adapop_workflow}) that conditions per-fact unlearning strength on an external popularity proxy and pairs this popularity scaling with an automated retain-balancing controller. AdaPop maps each forget example’s global popularity (e.g., inbound citelinks from Wikidata) to a per-fact exponent that reweights token ascent signals, then uses an epoch-level dual-ascent variable to enforce a bounded drift on a retain set. The novelty is two-fold: (1) explicit, externally-sourced popularity conditioning of per-fact gradient pressure, which directly addresses the popularity gap rather than proxying it via model confidence, and (2) an automated dual-ascent controller that operates on the popularity-weighted forget loss to find stable forget/retain operating points without dataset-specific tuning of loss weights. We empirically show AdaPop achieves deeper paraphrase-generalised erasure and larger internal representation disruption than competitive baselines while preserving general capabilities across multiple 7–8B LLM families and two real-world knowledge benchmarks, DUET~\cite{anna2026anatomy} and RWKU~\cite{jin2024rwku}.

\begin{comment}
To address this issue, we introduce \textbf{AdaPop} (\textbf{Ada}ptive \textbf{Pop}ularity), a popularity-weighted adaptive gradient unlearning method (Figure~\ref{fig:adapop_workflow}). AdaPop scales the gradient update for each forget-set fact according to its global popularity and automates the forget-retain balance via a dual-ascent controller. In a preprocessing stage, AdaPop uses Wikidata inbound citelink counts as a popularity annotator~\cite{anna2026anatomy} to assign each fact a per-fact exponent $\beta_i$ that controls the strength of its gradient update. Facts with a higher popularity score receive a stronger unlearning force; rare facts receive a weaker one.

We evaluate AdaPop against SOTA unlearning algorithms on two real-world knowledge benchmarks, DUET~\cite{anna2026anatomy} and RWKU~\cite{jin2024rwku}, across three model families (Llama-3.1-8B, Qwen2.5-7B, Gemma-7B). AdaPop consistently reduces forget-set ROUGE-L and cosine similarity below all baselines while keeping retain-set scores within the range of the unmodified model, with the largest gains on paraphrase-generalized and adversarial-robustness subsets, and with general language capabilities preserved throughout.

We further study the contribution of each AdaPop component through a controlled ablation (Appendix~\ref{app:ablation}), independently removing the popularity exponent $\beta_i$ and the dual-ascent controller. The results show that neither component alone accounts for the full gain: fixing $\alpha$ with dynamic $\beta_i$ is the worst-performing variant, while full AdaPop is the only configuration that produces both deep forget-set erasure and stable retain-set preservation across all models.

We hope AdaPop and our findings pave the way for future research on LLM unlearning via popularity salience.
\end{comment}

% \paragraph{Contributions.}
% \begin{itemize}
%     \item We confirm the \emph{popularity gap} as a systematic failure mode of confidence-based unlearning: methods calibrated on local model confidence under-erase popular facts while risking over-erasure of rare ones, a failure mode observed but not addressed by prior work~\cite{krishnannot, anna2026anatomy}.
%     \item We introduce \textbf{AdaPop}, which addresses this via (i) a popularity-dependent exponent $\beta_i$ derived from Wikidata to globally calibrate the unlearning signal per fact, and (ii) a dual-ascent controller that automatically adjusts the retain coefficient during training, eliminating manual hyperparameter search.
%     \item We evaluate AdaPop on DUET~\cite{anna2026anatomy} and RWKU~\cite{jin2024rwku} across three model families, showing that a global popularity signal produces deeper semantic erasure than local confidence-based methods across all models, with paraphrase-generalized forgetting as the clearest differentiator and general capabilities preserved throughout (§\ref{sec:main_results}--\ref{sec:lm_eval}).
% \end{itemize}

















% % VERSION 1
% Privacy regulations and the right to erasure create concrete demand for machine unlearning: removing specific knowledge from a trained model without retraining from scratch~\cite{mantelero2013eu, cao2015towards, bourtoule2021machine}. Large language models present a distinctive challenge. LLMs memorize substantial portions of their pretraining data~\cite{carlini2021extracting, tirumala2022memorization, carlini2023quantifying}, including personally identifiable information and sensitive facts, making post-hoc removal both practically necessary and technically difficult. The difficulty is not uniform: facts encountered frequently during pretraining are redundantly encoded across many layers and attention heads, making them far more resistant to gradient-based removal than rarely seen facts~\cite{kandpal2023large, hartmann2024undesirable, merullo2025linear}.

% Existing unlearning algorithms treat all facts in the forget set as equally hard to erase. Gradient Ascent (GA;~\citealt{jang2022knowledge}) maximizes the negative log-likelihood on the forget set without constraint, frequently causing catastrophic collapse of linguistic coherence~\cite{zhang2024negative}. Gradient Difference (GD;~\citealt{liu2022continual}) adds a retain loss but still assigns equal gradient weight to every forget token. Negative Preference Optimization (NPO;~\citealt{zhang2024negative}) reformulates unlearning as preference alignment, yielding a conservative signal that fails to erase well-memorized popular facts. Weighted Gradient Ascent (WGA;~\citealt{wang2025rethinking}) weights each token's gradient by the model's log-probability, concentrating unlearning on high-confidence tokens. All four methods calibrate the unlearning signal from the model's current output distribution, a quantity that changes throughout training and does not reflect how deeply a fact is parametrically encoded~\cite{hartmann2024undesirable, kandpal2023large}. The result is a systematic \emph{popularity gap}: methods calibrated for rare facts over-erase them while under-erasing popular ones, and vice versa.

% Krishnan et al.~\cite{krishnannot} confirm through controlled experiments across model scales that frequently encountered facts are either not unlearned at all or only superficially forgotten, remaining extractable via adversarial queries. Borisiuk et al.~\cite{anna2026anatomy} establish the same failure mode on the DUET benchmark, showing that fact salience systematically predicts unlearning difficulty across multiple gradient-based methods. Both works identify the problem empirically but offer no popularity-aware training method to address it.

% We introduce \textbf{AdaPop} (\textbf{Ada}ptive \textbf{Pop}ularity), which addresses this gap by conditioning gradient updates on an external popularity proxy and automating the forget-retain balance via a dual-ascent controller. As a popularity proxy, we use the sum of inbound Wikipedia citelinks to the subject and object nodes of each forget-set fact, computed from Wikidata (\texttt{pop\_sum}). This signal is static and independent of the model's current output distribution, providing stable per-fact calibration of unlearning strength. A dual-ascent controller adjusts the retain coefficient $\alpha$ at each training epoch to keep retain-set loss within a preset tolerance, eliminating manual hyperparameter search. Figure~\ref{fig:adapop_workflow} illustrates the full workflow.

% Through experiments on DUET~\cite{anna2026anatomy} and RWKU~\cite{jin2024rwku} across three model families (Llama-3.1-8B, Qwen2.5-7B, Gemma-7B), AdaPop produces deeper semantic erasure than all local confidence-based baselines, with the clearest gains on paraphrase-generalized forgetting metrics, while preserving general language capabilities across all settings.


% \begin{figure*}[t]
%     \centering
%     \includegraphics[width=0.85\textwidth]{latex/img/crop_adapop.pdf}
%     \caption{AdaPop assigns a per-fact exponent $\beta_i$ derived from an external popularity score, applying weaker unlearning to rare facts and stronger unlearning to popular ones. A dual-ascent controller adjusts the retain coefficient $\alpha$ at each epoch to keep retain-set loss within a preset tolerance.}
%     \label{fig:adapop_workflow}
% \end{figure*}


% \paragraph{Contributions.} We make the following contributions:
% \begin{itemize}
%     \item We confirm the \emph{popularity gap} as a systematic failure mode of confidence-based unlearning: methods calibrated on local model confidence under-erase popular facts while risking over-erasure of rare ones, a failure mode observed but not addressed by prior work~\cite{krishnannot, anna2026anatomy}.
%     \item We introduce \textbf{AdaPop}, which addresses this via (i) a popularity-dependent exponent $\beta_i$ derived from Wikidata to globally calibrate the unlearning signal per fact, and (ii) a dual-ascent controller that automatically adjusts the retain coefficient during training, eliminating manual hyperparameter search.
%     \item We evaluate AdaPop on DUET~\cite{anna2026anatomy} and RWKU~\cite{jin2024rwku} across three model families (Llama-3.1-8B, Qwen2.5-7B, Gemma-7B), showing that a global popularity signal produces deeper semantic erasure than local confidence-based methods across all models, with paraphrase-generalized forgetting as the clearest differentiator and general capabilities preserved throughout (§\ref{sec:main_results}--\ref{sec:lm_eval}).
% \end{itemize}









%% PLAN FROM 14/03/26
% %% Describe what is the problem current algos (stability isses, popularity influence): Gradient Ascent, Gradient Difference, NPO, WGA;
% %% Explain how we solve this popularity gap 
% %% Which benchmark we use to evaluate (we use real-world facts QA benchmarks for solving real-words tasks; we scoring facts based on they popularity. if you have ascess to pretrain set, you can use information from it to scoring populaeiry; otherwise use 'approximate' metrics from large data corpus like wikidata to score facts. here popularity is the global metric for for unlerning regularization; the local (model) metric is log probes and its to factors influence to unlearning.) also for trade-off stability we automatically balance retain, that directly effect to forgetting.

% %% добавить ссылки, можно посмотреть из интро статьи (https://arxiv.org/pdf/2504.05058)
% % Privacy regulations and the right to erasure create concrete demand for machine unlearning: the ability to remove specific facts from a trained model without retraining from scratch. In LLMs, however, removing a fact is not symmetric with inserting it.

% Cуществующие сота анлеарн методы ортимизируют лоссы типа преференс оптиизации \cite{NPO} и градиентной разницы \cite{GA, GD, WGA}.
% Эти methods rely on local model confidence as their primary signal, and none adjust the forget-retain balance dynamically. 
% Однако,не в токенам дело
% Таким образом, 
% проблемы методов на реал ворлд дата в том, что все точки одинаковые krishnannot и мы. Однако известно, что частота фактов влияет на меморизацию ссылка \cite{}.  

% To address the issue, мы предлагаем метод адаптивного анлергинга на популярности named adapop. As shown in Figure,. We introduce \textbf{AdaPop} (Adaptive Popularity), which addresses this gap by scaling gradient updates per fact using an external popularity proxy and automates the forget-retain balance via a dual-ascent controller. We introduce \textbf{AdaPop}, which addresses this via (i) a popularity-dependent exponent $\beta_i$ derived from Wikidata globally calibrate the unlearning signal per fact, and (ii) a dual-ascent controller that automatically adjusts the retain coefficient during training, eliminating manual hyperparameter search.


% Through extensive experiments across two benchmarks such as DUET and RWKU, we demonstrate constant improvement on forget set while preserving retain results. В частности, цифраю 

% %Facts encountered frequently during pretraining are redundantly encoded across many layers and attention heads, making them far more resistant to gradient-based removal than rarely encountered facts. Existing algorithms treat all facts in the forget set as equally hard to erase, producing a systematic popularity gap: methods calibrated for rare facts over-erase them while under-erasing popular ones, and methods calibrated for popular facts risk collateral damage to rare knowledge and general capabilities \cite{anna2026anatomy, krishnannot}.

% %Gradient Ascent (GA, \cite{jang2022knowledge}) maximizes loss on the forget set without constraint, which frequently leads to catastrophic collapse of linguistic coherence. Gradient Difference (GD, \cite{liu2022continual}) adds a retain loss but still treats every forget token as equally important. Negative Preference Optimization (NPO, \cite{zhang2024negative}) reformulates unlearning as preference alignment, yielding a conservative signal that fails to erase well-memorised popular facts. Weighted Gradient Ascent (WGA, \cite{wang2025rethinking}) weights gradient updates by the model's own log-probabilities, focusing ascent on high-confidence tokens. All above-mentioned methods rely on local model confidence as their primary signal, and none adjust the forget-retain balance dynamically. The result is a suboptimal tradeoff: methods must be hand-tuned per dataset and model to avoid either insufficient forgetting or retain-set degradation. Our central claim is that log-probabilities alone do not accurately reflect how deeply a fact is embedded in the model~\cite{hartmann2024undesirable, kandpal2023large}; an external global popularity proxy provides a more reliable signal for calibrating unlearning strength.

% \begin{figure}[ht!]
%     \centering
%     \includegraphics[width=1\linewidth]{latex/img/ada_pop_sq_illustr.jpg}
%     \caption{The proposed AdaPop assigns a per-fact exponent $\beta_i$ from an external popularity score, applying weaker unlearning to rare facts and stronger unlearning to popular ones. A dual-ascent controller adjusts the retain coefficient $\alpha$ at each epoch to keep retain-set loss within a preset tolerance.}
%     \label{fig:adapop_workflow}
% \end{figure}


%  As a popularity proxy, we use the sum of inbound citelinks (Wikipedia page references) to an entity's subject and object nodes, computed directly from Wikidata for each fact in the forget set; this quantity is Wikidata \texttt{pop\_sum}.

% \paragraph{Our contributions:} %We make the following contributions:
% % (Figure~\ref{fig:adapop_workflow})
% \begin{itemize}
%     \item We confirm the \emph{popularity gap} as a systematic failure mode of confidence-based unlearning: methods that calibrate updates on local model confidence under-erase popular facts (high memory inertia) while risking over-erasure of rare ones.

%     \item  We design a method for factuality-aware MU that trains something на частоте.    %\item We introduce \textbf{AdaPop}, which addresses this via (i) a popularity-dependent exponent $\beta_i$ derived from external knowledge bases (e.g., Wikidata) to globally calibrate the unlearning signal per fact, and (ii) a dual-ascent controller that automatically adjusts the retain coefficient during training, eliminating manual hyperparameter search.
    
%     %\item We evaluate \textbf{AdaPop} on DUET~\cite{anna2026anatomy} and RWKU~\cite{jin2024rwku} datasets across three models from different families (Llama-3.1-8B, Qwen2.5-7B, Gemma-7B), showing that a global popularity signal produces deeper semantic erasure than local confidence-based methods across all models, with paraphrase-generalised forgetting as the clearest differentiator and general capability preserved throughout (§\ref{sec:main_results}--\ref{sec:lm_eval}).
%     \item We demonstrate that on two benchmark datasets, our method outperforms state-of-the-art on both MU and erase/attack metrics.

% \end{itemize}



\section{Related Work}
\label{sec:related}
% % \subsection{Gradient-Based Machine Unlearning}
% Weight-editing methods erase target facts by modifying model parameters directly, motivated by privacy regulations and the right to erasure~\cite{yao2024large}. Early approaches relied on prompting-based or in-context interventions~\cite{liu2024survey}; Łucki et al.~\cite{lucki2025adversarial} show that such inference-time methods leave parametric knowledge intact, with models remaining vulnerable to adversarial extraction under targeted queries. Parameter-editing methods address this gap but introduce a different challenge: none accounts for variation in how deeply individual facts are encoded across the model.

% Gradient Ascent (GA) maximizes the negative log-likelihood on the forget set, which frequently leads to catastrophic collapse of linguistic coherence~\cite{yao2024large, zhang2024negative}. Gradient Difference (GD) adds a retain loss to bound collateral damage but treats every token in the forget set with equal importance~\cite{bu2025unlearning}. Negative Preference Optimization (NPO) frames unlearning as preference alignment, pushing the output distribution away from the memorised response~\cite{zhang2024negative}. Weighted Gradient Ascent (WGA) weights each token's gradient by the model's log-probability, concentrating unlearning effort on high-confidence tokens~\cite{wang2025rethinking}. All four methods calibrate the unlearning signal from the model's current output distribution, which changes throughout training and does not reflect how deeply a fact is parametrically encoded~\cite{zhang2024negative, wang2025rethinking}. Methods that apply a retain loss (GD, WGA) further require manual selection of the retain coefficient $\alpha$: too small and general utility degrades; too large and target facts survive erasure. AdaPop replaces this local signal with a global popularity proxy and eliminates manual retain tuning with a dual-ascent controller that adjusts $\alpha$ automatically, directly addressing both the per-fact calibration gap and the hyperparameter sensitivity that none of these methods resolves.

% \subsection{Knowledge Popularity and Memorization}
% Kandpal et al.~\cite{kandpal2023large} show that LLMs recall long-tail facts far less reliably than frequent ones, and Hartmann et al.~\cite{hartmann2024undesirable} survey how parametric memorisation scales with training-corpus frequency. Popular facts are encoded redundantly across many layers; rare facts are lightly encoded and susceptible to over-erasure when subjected to the same gradient force. This asymmetry means a single gradient weight cannot address both regimes simultaneously.

% Log-probabilities do not accurately reflect memorisation depth~\cite{zhang2024negative}: a model may assign high confidence to a surface token pattern without the corresponding fact being deeply encoded across layers. External popularity proxies, such as entity citelink counts or co-occurrence frequencies, provide a fact-level signal independent of the model's current output distribution {\cite{krishnannot, anna2026anatomy}}. AdaPop assigns each fact a popularity-dependent exponent, treating memorisation depth as a property of the fact rather than of the model.


% VERSION 3
\subsection{Machine Unlearning Methods}
\label{subsec:gradient_unlearning}

Gradient Ascent (GA;~\citealt{jang2022knowledge}) maximizes the negative log-likelihood on the forget set without constraint, frequently causing catastrophic collapse of linguistic coherence~\cite{yao2024large, zhang2024negative}. Gradient Difference (GD;~\citealt{liu2022continual}) bounds collateral damage with a retain loss but assigns equal gradient weight to every forget token~\cite{bu2025unlearning}. Negative Preference Optimization (NPO;~\citealt{zhang2024negative}) frames unlearning as preference alignment, pushing the output distribution away from the memorized response. Weighted Gradient Ascent (WGA;~\citealt{wang2025rethinking}) weights each token's gradient by the model's log-probability, concentrating unlearning on high-confidence tokens. All four calibrate the unlearning signal from the model's current output distribution, which changes throughout training and does not reflect the degree to which a fact is parametrically encoded. \citet{lucki2025adversarial} further show that inference-time interventions leave parametric knowledge intact, with models remaining vulnerable to adversarial extraction under targeted queries.

\subsection{Memorization, Frequency, and Unlearning}
\label{subsec:memorization}

The extent of LLM memorization scales with training-corpus frequency~\cite{carlini2023quantifying, tirumala2022memorization}: \citet{kandpal2023large} show LLMs recall long-tail facts far less reliably than frequent ones, and \citet{merullo2025linear} demonstrate that frequent facts are encoded more linearly in model representations, providing a mechanistic basis for why popular facts resist gradient-based removal. This asymmetry is adversarial to unlearning: verbatim memorization impedes removal in both vision~\cite{NEURIPS2024_16e18fa3} and LLM settings~\cite{barbulescu2024each}, and \citet{baluta2024unlearning} confirm in a controlled setup that differently exposed data points are unlearned differently. Log-probabilities do not accurately reflect memorization depth~\cite{wang2025rethinking, hartmann2024undesirable}: a model may assign high confidence to a surface pattern even when the corresponding fact is not deeply encoded. External popularity proxies, such as entity citelink counts, provide a stable fact-level signal independent of the model's output distribution.

\subsection{Popularity Effects on Unlearning}
\label{subsec:popularity_effects}

\citet{krishnannot} establish through controlled pretraining experiments that high-frequency facts are either not unlearned or only superficially forgotten, remaining extractable via adversarial queries across multiple methods and model scales. \citet{anna2026anatomy} confirm the same failure mode on real-world facts in the DUET benchmark, showing that Wikipedia-based fact salience systematically predicts unlearning difficulty across SOTA unlearning methods. Both works precisely characterize the popularity gap and provide benchmarks we build on, but neither proposes a training method to address it. AdaPop is, to our knowledge, the first unlearning method that conditions per-fact gradient updates on an external popularity signal.





%% PLAN FROM 14/03/26
% \subsection{Gradient-Based Machine Unlearning}
% Weight-editing methods erase target facts by directly modifying model parameters, motivated by privacy regulations and the right to erasure~\cite{yao2024large}. Early approaches relied on prompting-based or in-context interventions~\cite{liu2024survey}; \citet{lucki2025adversarial} show that such inference-time methods leave parametric knowledge intact, with models remaining vulnerable to adversarial extraction under targeted queries. Parameter-editing methods address this gap but introduce a different challenge: none accounts for variation in how deeply individual facts are encoded across the model.

% Gradient Ascent (GA) maximizes the negative log-likelihood on the forget set, which frequently leads to catastrophic collapse of linguistic coherence~\cite{yao2024large, zhang2024negative}. Gradient Difference (GD) adds a retain loss to bound collateral damage but treats every token in the forget set with equal importance~\cite{bu2025unlearning}. Negative Preference Optimization (NPO) frames unlearning as preference alignment, pushing the output distribution away from the memorised response~\cite{zhang2024negative}. Weighted Gradient Ascent (WGA) weights each token's gradient by the model's log-probability, concentrating unlearning effort on high-confidence tokens~\cite{wang2025rethinking}. All four methods calibrate the unlearning signal from the model's current output distribution, which changes throughout training and does not reflect how deeply a fact is parametrically encoded~\cite{zhang2024negative, wang2025rethinking}. AdaPop replaces this local signal with a global popularity proxy, directly addressing the per-fact calibration gap that none of these methods resolves.

% \subsection{Knowledge Popularity and Memorization}
% \citet{kandpal2023large} show that LLMs recall long-tail facts far less reliably than frequent ones, and \citet{hartmann2024undesirable} survey how parametric memorisation scales with training-corpus frequency. Popular facts are redundantly encoded across many layers; rare facts are lightly encoded and susceptible to over-erasure when subjected to the same gradient. This asymmetry means a single gradient weight cannot address both regimes simultaneously.

% наиболее близкие к нам работы в MU - {\cite{krishnannot, anna2026anatomy}}. Они оценили сота методы такие как Х1 Х2 Х3 и экспериментально подвердили фандинг проблем. Однако, не предлагали свое решение. To the best of our knowledge, our proposed method is the first pop-aware MU method to fix the issue. 

% Log-probabilities do not accurately reflect memorisation depth~\cite{zhang2024negative}: a model may assign high confidence to a surface token pattern even when the corresponding fact is not deeply encoded across layers. External popularity proxies, such as entity citelink counts or co-occurrence frequencies, provide a fact-level signal independent of the model's current output distribution {\cite{krishnannot, anna2026anatomy}}. 

% %AdaPop assigns each fact a popularity-dependent exponent, treating memorisation depth as a property of the fact rather than of the model.


\section{Adaptive Popularity Method (AdaPop)}
\label{sec:AdaPop}
%AdaPop is a gradient-based weight-editing method: this family directly modifies parametric memory at training time, scales to arbitrary knowledge types without architectural assumptions, and is the class in which the popularity gap-differential resistance of facts to equal gradient force-manifests most directly. Methods such as GD, NPO, and WGA require manual tuning of unlearning hyperparameters like the retain coefficient $\alpha$: too small and the model loses general utility; too large and target facts survive erasure. AdaPop replaces this manual search with two components: a popularity-dependent exponent $\beta_i$ that calibrates the per-fact unlearning signal from an external metric, and a dual-ascent controller that adjusts $\alpha$ automatically during training.
We frame targeted unlearning as a constrained optimisation: increase the forget-set loss while bounding the retain-set drift. The proposed AdaPop solves this by combining (i) per-fact popularity scaling via a power-law exponent computed from an external popularity proxy that reweights per-token ascent signals, and (ii) a dual-ascent retain controller that updates the retain penalty at epoch granularity. Together, these components close the ``popularity gap'' by (a) concentrating the ascent signal on the highest-confidence answer tokens for rare facts, where encoding is narrow and over-erasure risk is high, and (b) distributing the ascent signal uniformly across all answer tokens for popular facts, where broad parametric encoding demands wide gradient coverage.



\subsection{The Popularity Gap}

% TODO: It should be noted that log-probability does not accurately reflect the actual popularity of a fact, as it is subject to noise from prepositions and particles, which typically have a high probability.
Parametric memorisation scales with training-corpus frequency~\cite{hartmann2024undesirable, kandpal2023large}. A single shared $\alpha$ (as in GD) or a confidence-based weight (as in WGA) cannot simultaneously address both regimes: popular facts resist erasure while rare facts risk collateral damage from the same gradient magnitude. AdaPop assigns each fact an independent exponent $\beta_i$ derived from an external popularity score, so that the unlearning signal is calibrated to fact-level memorisation depth rather than to the model's current output distribution.



\subsection{Problem Formulation}

Let $\mathcal{F}=\{(x_i,y_i,s_i)\}_{i=1}^{N_F}$ be the \emph{forget} set, where each entry consists of a question-answer pair $(x_i, y_i)$ and a corresponding popularity score $s_i$ derived from a global metric such as Wikidata citelinks count (\texttt{pop\_sum}). Let $\mathcal{R}=\{(x_j,y_j)\}_{j=1}^{N_R}$ be the \emph{retain} set. We consider an autoregressive model $q_\theta$ defined by the conditional probability $q_\theta(y_t \mid x, y_{<t})$ of a token $y_t$ given its prefix. For a token position $t$ in sequence $i$, we define the negative log-likelihood (NLL)
\begin{equation}
    \ell_{i,t}(\theta) = -\log q_\theta(y_{i,t} \mid x_i, y_{i,<t})
\end{equation}
and the model's confidence
\begin{equation}
    p_{i,t}(\theta) = \exp(-\ell_{i,t}(\theta)).
\end{equation}

The sets $\Omega_{\mathcal{F}}$ and $\Omega_{\mathcal{R}}$ represent the indices of valid tokens not masked by padding or special tokens. The objective of AdaPop is to increase the loss on $\mathcal{F}$ while bounding the drift on $\mathcal{R}$. We frame this as a constrained optimisation problem (maximise forgetting subject to $\delta_k \le \varepsilon$) and enforce it via a dual-ascent controller (§\ref{sec:dual_ascent}). A fixed retain coefficient $\alpha$, as in GD~\cite{liu2022continual} and WGA~\cite{wang2025rethinking}, requires dataset-specific grid search; the constrained formulation eliminates this by adjusting $\alpha$ automatically in response to observed retain drift. Full design rationale is in Appendix~\ref{app:design_motivation}.



\subsection{Popularity-Aware Forget Objective}

Each sample $i$ in the forget set is assigned an exponent $\beta_i$ based on its popularity score $s_i$:
\begin{equation}
    \beta_i = a \cdot s_i^{-b},
\end{equation}
A power-law mapping is natural here because parametric memorisation itself scales as a power law with corpus frequency~\cite{kandpal2023large}: since $s_i$ proxies that frequency, $\beta_i \propto s_i^{-b}$ mirrors the underlying relationship. This choice of power law compresses the proxy’s wide dynamic range into a bounded exponent that places rare facts in a self-limiting regime ($\beta>1$) and popular facts in a pressure-sustaining regime ($\beta<1$), improving forget/retain tradeoffs. The coefficients $a$ and $b$ are proxy-specific (see Appendix~\ref{app:derivation} for derivation). $\beta(s)$ is monotonically decreasing: as popularity increases, $\beta_i$ decreases, distributing the ascent signal more broadly across tokens for popular facts. For rare facts, a larger $\beta_i$ concentrates the update on the highest-confidence tokens, limiting collateral drift. Values are clipped to $[0.05, 2.0]$ in all experiments to avoid numerical instability.

The per-token unlearning weight is:
\begin{equation}
w_{i,t} = \left(\mathrm{stopgrad}(p_{i,t}(\theta))\right)^{\beta_i}.
\end{equation}
The forget loss is:
\begin{equation}
L_f(\theta) = -\frac{1}{|\Omega_{\mathcal{F}}|} \sum_{(i,t) \in \Omega_{\mathcal{F}}} w_{i,t} \ell_{i,t}(\theta).
\end{equation}

%\subsection{Implementation Note}

The weighting $w_{i,t}$ is applied only during training, not at inference. It is computed only over answer tokens $y_i$; question tokens $x_i$ are masked and do not contribute to the loss, so the model's ability to parse the query remains unaffected.





\subsection{Dual-Ascent and Auto-Balancing}\label{sec:dual_ascent}

AdaPop treats $\alpha$ as a dual variable enforcing the constraint $\delta_k \le \varepsilon$, where $\delta_k$ is the one-sided relative drift of the retain loss $R_k$ at epoch $k$:
\begin{equation}
\delta_k = \max\left( 0, \frac{R_k - R_{\mathrm{ref}}}{\max(R_{\mathrm{ref}}, \epsilon_0)} \right),
\end{equation}
$R_{\mathrm{ref}}$ is the retain loss measured at the end of the first epoch, and $\epsilon_0 > 0$ is a small constant for numerical stability.

\paragraph{Epoch-Level Dual Update.}
Let $\alpha_k = \alpha_0 + \lambda_k$. The dual variable $\lambda$ is updated via projected gradient ascent once per epoch (not per step) to reduce oscillation due to noisy per-batch retain losses:
\begin{align}
    \lambda_{k+1} &= \Pi_{[0, \lambda_{\max}]}\left( \lambda_k + \eta_\lambda (\delta_k - \varepsilon) \right) \\
    \alpha_{k+1} &= \Pi_{[0,\, \alpha_0 + \lambda_{\max}]}(\alpha_0 + \lambda_{k+1}),
\end{align}
where the dual updates enforce the retain constraint on the popularity-weighted forget loss (rather than the unweighted or entropic formulations), aligning the feedback signal with the actual gradient mass distribution and enabling the controller to find operating points with deeper, targeted forgetting. We empirically found both choices important for stability (see Fig.~\ref{fig:ablation} in the Appendix).

\paragraph{Per-Step Popularity Modulation (optional).}
AdaPop supports an optional fine-grained control mode that scales the effective retain coefficient per batch. Since batches containing popular facts (small $\beta$) generate aggressive unlearning signals, the effective retain coefficient $\alpha_{\mathrm{eff}} = \alpha \cdot \text{scale}$ can be boosted to compensate:
\begin{equation}
\text{scale} = \mathrm{clip}\!\left(\frac{1}{1 + g \cdot \bar{\beta}},\; s_{\min},\; s_{\max}\right)
\end{equation}
where $g \geq 0$ is a gain coefficient, $\bar{\beta}$ is the batch-mean exponent, and $[s_{\min}, s_{\max}]$ clips the scale to a stable range. When $g = 0$, $\text{scale} = 1$, and the batch-level modulation is disabled. In all experiments reported in this paper, $g = 0$ (i.e., per-step modulation is off); the dual-ascent epoch-level controller alone provides the retain's balance.

\begin{algorithm}[t]
\caption{AdaPop}
\label{alg:adapop}
\begin{algorithmic}[1]
\REQUIRE Forget set $\mathcal{F}=\{(x_i,y_i,s_i)\}$ with popularity scores $s_i$, retain set $\mathcal{R}$, model $q_\theta$, initial retain weight $\alpha_0$, dual step size $\eta_\lambda$, tolerance $\varepsilon$, anchor points $(s_r, s_p)$
\ENSURE Updated parameters $\theta$

\STATE Compute popularity exponents:
\(
\beta_i = \mathrm{clip}(a\, s_i^{-b}, \beta_{\min}, \beta_{\max})
\)
where coefficients $a,b$ are derived from anchor constraints $\beta(s_r)$ and $\beta(s_p)$.

\STATE Initialize $\lambda \leftarrow 0$, $\alpha \leftarrow \alpha_0$

\FOR{epoch $k = 1$ to $E$}

    \FOR{batch $B$ sampled from $\mathcal{F} \cup \mathcal{R}$}

        \FOR{each forget example $i$ and answer token $t$ in $B$}
            \STATE $p_{i,t} \leftarrow \mathrm{stopgrad}(\exp(-\mathrm{NLL}_{i,t}(\theta)))$
            \STATE $w_{i,t} \leftarrow p_{i,t}^{\beta_i}$
        \ENDFOR

        \STATE Compute weighted forget loss:
        \[
        L_f = - \frac{1}{|B_F|} 
        \sum_{(i,t)\in B_F} w_{i,t}\,\mathrm{NLL}_{i,t}(\theta)
        \]

        \STATE Compute retain loss $L_r$ over retain tokens

        \STATE Total loss:
        \(
        L = L_f + \alpha L_r
        \)

        \STATE Update model parameters $\theta$ using gradient descent (LoRA updates)

    \ENDFOR

    \STATE Compute retain drift:
    \(
    \delta_k = \frac{\max(0, R_k - R_{\text{ref}})}{\max(R_{\text{ref}}, \epsilon_0)}
    \)

    \STATE Dual update:
    \(
    \lambda \leftarrow \mathrm{proj}_{[0,\lambda_{\max}]}\big(\lambda + \eta_\lambda(\delta_k - \varepsilon)\big)
    \)

    \STATE Update retain weight:
    \(
    \alpha \leftarrow \mathrm{proj}_{[0,\alpha_0+\lambda_{\max}]}(\alpha_0 + \lambda)
    \)

\ENDFOR

\RETURN $\theta$
\end{algorithmic}
\end{algorithm}

Two design choices in the proposed AdaPop (Alg.~\ref{alg:adapop}) are worth highlighting. The dual update runs at epoch granularity: per-batch retain losses are noisy, and epoch-level feedback reduces oscillation while still responding to retain drift within a few hundred gradient steps. The controller uses a popularity-weighted forget loss rather than an unweighted formulation, aligning the feedback signal with the actual gradient mass distribution and enabling the controller to find operating points with deeper, targeted forgetting. Ablations in Appendix~\ref{app:ablation} confirm both choices are necessary.



%% take from other/ada_pop.tex

% In this section, we detail the \textbf{AdaPop} (Adaptive Popularity) framework. Our approach moves beyond local model confidence by integrating a global popularity metric and an automated dual-ascent controller to maintain the balance between knowledge erasure and model utility.

% \subsection{Motivation and Setting}

% Let $\mathcal{F}=\{(x_i,y_i,s_i)\}_{i=1}^{N_F}$ be the \emph{forget} set, where each entry consists of a question-answer pair $(x_i, y_i)$ and a corresponding popularity score $s_i$ derived from a global metric such as \texttt{pop\_sum}. Let $\mathcal{R}=\{(x_j,y_j)\}_{j=1}^{N_R}$ be the \emph{retain} set used to preserve general model performance. We consider an autoregressive model $q_\theta$ where the conditional probability of a token $y_t$ given its prefix is $q_\theta(y_t \mid x, y_{<t})$. For any token position $t$ in sequence $i$, we define the negative log-likelihood (NLL) and the model's confidence as:
% \begin{align}
%     \ell_{i,t}(\theta) &= -\log q_\theta(y_{i,t} \mid x_i, y_{i,<t}) \\
%     p_{i,t}(\theta) &= \exp(-\ell_{i,t}(\theta))
% \end{align}
% The sets $\Omega_{\mathcal{F}}$ and $\Omega_{\mathcal{R}}$ represent the indices of valid tokens not masked by padding or special tokens. The objective of AdaPop is to increase the loss on $\mathcal{F}$ while minimizing the drift on $\mathcal{R}$ through a weighted gradient ascent objective where the retain coefficient $\alpha$ is dynamically adjusted.

% \subsection{Popularity-Aware Forget Objective}

% A core contribution of AdaPop is the use of a \textbf{popularity-dependent exponent} $\beta_i$ to scale the unlearning signal. Each sample $i$ in the forget set is assigned an exponent based on the popularity score $s_i$ of the corresponding fact. The exponent is defined as a general function of popularity:
% \[
% \beta_i = f(s_i).
% \]
% In practice, this function is often defined empirically to control the strength of the update. For example, one possible form used in our implementation is:
% \[
% \beta_i = 58.7 s_i^{-0.796},
% \]
% where the constant $58.7$ and the exponent $-0.796$ were chosen empirically to balance the influence across the popularity spectrum. To ensure training stability and prevent gradient explosion, $\beta_i$ is typically clipped to a stable interval (e.g., $[0.05, 2.0]$). Alternatively, a constant exponent $\beta_i = \beta_{\mathrm{const}}$ may be used to treat all facts with equal importance.



% The relationship between popularity and the final unlearning weight $w_{i,t}$ is governed by the detached token weight:
% \[
% w_{i,t} = \left(\mathrm{stopgrad}(p_{i,t}(\theta))\right)^{\beta_i}.
% \]
% This formulation creates a critical inverse relationship. Popular facts (high $s_i$) result in a smaller $\beta_i$, which amplifies and "flattens" the unlearning signal. This ensures that even tokens with moderate confidence contribute significantly to the gradient ascent, overcoming the strong "memory inertia" of deeply embedded knowledge. Conversely, rare facts (low $s_i$) receive a larger exponent, which "sharpens" the update to concentrate only on the most certain tokens, preventing unnecessary drift on niche information. The final forget loss is defined as:
% \[
% L_f(\theta) = -\frac{1}{|\Omega_{\mathcal{F}}|} \sum_{(i,t) \in \Omega_{\mathcal{F}}} w_{i,t} \ell_{i,t}(\theta).
% \]

% \subsection{Dual-Ascent and Fine-Grained Retain Control}

% To maintain stability without manual hyperparameter tuning, AdaPop employs a two-tier control system for the retain coefficient $\alpha$.

% \paragraph{Epoch-Level Dual Update.} 
% AdaPop treats $\alpha$ as a dual variable to satisfy the constraint $\delta_k \le \varepsilon$, where $\delta_k$ is the one-sided relative drift of the retain loss $R_k$ at epoch $k$:
% \[
% \delta_k = \max\left( 0, \frac{R_k - R_{\mathrm{ref}}}{\max(R_{\mathrm{ref}}, \epsilon_0)} \right).
% \]
% Let $\alpha_k = \alpha_0 + \lambda_k$. The dual variable $\lambda$ is updated via projected gradient ascent:
% \begin{align}
%     \lambda_{k+1} &= \Pi_{[0, \lambda_{\max}]}\left( \lambda_k + \eta_\lambda (\delta_k - \varepsilon) \right) \\
%     \alpha_{k+1} &= \Pi_{[0,\, \alpha_0 + \lambda_{\max}]}(\alpha_0 + \lambda_{k+1})
% \end{align}

% \paragraph{Per-Step Popularity Modulation.} 
% Beyond epoch-level updates, AdaPop employs fine-grained control to protect the model from sudden collapse. Since batches containing popular facts (small $\beta$) generate aggressive unlearning signals, we proactively scale the effective retain coefficient $\alpha_{\mathrm{eff}}$ for each batch:
% \[
% \text{scale} = \frac{1}{1 + \text{gain} \cdot \bar{\beta}}
% \]
% This scaling strengthens the "shield" of the retain loss in real-time when the model is performing the loud, wide updates required for popular facts, ensuring that the unlearning process does not cause immediate collateral damage to general knowledge.

\section{Experimental Setup}
\label{sec:experiments}
% %% duet, rwku
% %% metrcis description: rouge-l recall, cos_sim + more deep interp on best lr, + std
% %% explain why we dont use something like tofu (its syntetic benchmark and there's only 20 QA to each pearson that is very small for popularity)


%% Explain why we dont use algo like RMU, prompt tuning and rejection tuning and etc
%% Describe benchmarks that i use: DUET (more diversety popularity), RWKU (model know knowlage from this benchmark perfectly, so we dont tune model on it) 


% %% explain that on datasets when popularity divercity is small (all dataset have simular pop distribution), we dont see so strong effect. but on duet, we can see the difference.

% %% explain approximation of dataset if we dont have eccess to full pretrain courpus by large coupose like wikidata/ commonrow 


%% Sections 4–5: Experimental Setup and Results

We evaluate \textbf{AdaPop} against SOTA unlearning methods: GA~\cite{jang2022knowledge}, GD~\cite{liu2022continual}, NPO~\cite{zhang2024negative}, and WGA~\cite{wang2025rethinking} across three LLM families and two real-world knowledge benchmarks. All experiments use LoRA fine-tuning~\cite{hu2021lora} with rank~$r{=}32$, $\alpha{=}64$, and a fixed batch size of 1 with 32 gradient accumulation steps. Main results are reported at learning rate $10^{-4}$, which provides a stable operating point where all methods have non-zero gradients while avoiding the catastrophic collapse that uniformly afflicts unconstrained ascent methods at higher rates. A full sweep over five learning rates is reported in Appendix~\ref{app:lr_sensitivity}. AdaPop-specific parameters: initial retain coefficient $\alpha_0 = 0.5$, constraint threshold $\varepsilon = 0.1$, dual step size $\eta_\lambda = 0.1$, dual variable bound $\lambda_{\max} = 5.0$, popularity exponent coefficients $a{=}58.7$, $b{=}0.796$, clipped to $[0.05, 2.0]$. Per-step popularity modulation is disabled ($g=0$). Sensitivity of results to $\pm$20\% perturbations of $a$ and $b$ is reported in Appendix~\ref{app:coef_sensitivity}.

\subsection{Models}
We use Llama-3.1-8B-Instruct~\cite{grattafiori2024llama}, Gemma-7B-it~\cite{team2024gemma}, and Qwen2.5-7B-Instruct~\cite{qwen2025}. The 7--8B scale is standard in unlearning research~\cite{zhang2024negative, yao2024large, jin2024rwku}. These models share the transformer backbone but differ in pretraining and instruction-tuning regimes; per-model breakdowns appear in Appendix~\ref{app:model_specific}.

\subsection{Benchmarks}
We evaluate on two real-world knowledge datasets that exhibit fundamentally different popularity distributions (Figure~\ref{fig:benchmarks_pop_diverse}):
\begin{itemize}
    \item \textbf{DUET}~\cite{anna2026anatomy}: A city-centric factual QA benchmark with explicit popularity diversity. It covers a broad spectrum of Wikidata \texttt{pop\_sum} scores (range: 69--3{,}763; median: 1{,}089.5), enabling direct observation of the popularity gap between the rare and popular fact subsets. We use the combined rare-and-popular forget subset alongside a 500-sample fast retain set.
    \item \textbf{RWKU}~\cite{jin2024rwku}: A precision-focused benchmark testing targeted fact erasure while preserving semantically adjacent knowledge. Its popularity distribution is narrower and skewed to the lower end (range: 0--704; median: 130), making it suitable for evaluating whether the accuracy of forgetting degrades under lower popularity variance.
\end{itemize}

\begin{figure}[t]
    \centering
    \includegraphics[width=0.9\linewidth]{latex/img/pop_sum_diversity_comparison.png}
    \caption{Wikidata \texttt{pop\_sum} score distributions for the forget sets of DUET and RWKU. DUET spans nearly two orders of magnitude ($69$--$3{,}763$; median $1{,}090$), providing a direct test of popularity-sensitive methods; RWKU is narrower and skewed lower (median $130$).}
    \label{fig:benchmarks_pop_diverse}
\end{figure}

We exclude synthetic benchmarks such as TOFU~\cite{maini2024tofu}, whose narrow popularity distributions are insufficient to evaluate popularity-sensitive methods. DUET provides \texttt{pop\_sum} natively; for RWKU, we compute it by querying Wikidata for the inbound citelink counts of each fact's subject and object entities. Because widely-referenced entities appear proportionally across all major web-scale pretraining corpora, Wikidata citelink counts correlate with training-corpus frequency regardless of the specific pretraining regime~\cite{kandpal2023large}, making \texttt{pop\_sum} a corpus-agnostic proxy applicable across the three model families we evaluate. LLM-estimated knowledgeability scores are a practical alternative when Wikidata access is unavailable~\cite{anna2026anatomy, marina2025llm}.


% ─────────────────────────────────────────────────────────────────────────────
\subsection{Evaluation Metrics}

We report six complementary metrics that jointly measure erasure depth and model integrity, grouped into output-level (§\ref{sec:main_results}) and internal representation metrics (§\ref{sec:internal}).

\begin{itemize}
    \item \textbf{ROUGE-L Recall}: token overlap between generated and gold answer. $\downarrow$ on forget, $\uparrow$ on retain.
    \item \textbf{Cosine Similarity (Cos Sim)}: embedding-space proximity to the gold label. Same directions.
    \item \textbf{$\Delta$ Log-Probability ($\Delta$LP)}: per-token log-likelihood shift from the base model. Large negative on forget; near-zero on retain.
    \item \textbf{$\Delta$ Rank}: change in ground-truth token rank in the predicted distribution. Large positive on forget; near-zero on retain.
    \item \textbf{Hidden-State Cosine Similarity (Hid.Cos)}: mean cosine of intermediate hidden states vs.\ the base model. Low on forget; high on retain.
    \item \textbf{KL Divergence (KL)}: KL of the output distribution from the base model. High on forget; low on retain.
\end{itemize}

% ─────────────────────────────────────────────────────────────────────────────
\section{Results}
\label{sec:results}

We structure the evaluation in four parts: output-level forgetting and retention (§\ref{sec:main_results}), robustness to paraphrase and adversarial queries (§\ref{sec:robustness}), internal representation analysis (§\ref{sec:internal}), and general capability preservation (§\ref{sec:lm_eval}).

\subsection{Main Results: ROUGE-L and Cosine Similarity}
\label{sec:main_results}

Tables~\ref{tab:rouge-llama}--\ref{tab:rouge-gemma} report ROUGE-L Recall and Cosine Similarity on the forget (F.) and retain (R.) splits at $\text{lr}=10^{-4}$. ``Orig'' refers to the unmodified pre-LoRA base model. Across all three architectures and both benchmarks, \textbf{AdaPop} achieves the lowest forget ROUGE-L among stable methods while keeping retain ROUGE-L above $0.855$.

% --- Table 1: Llama-3.1-8B-Instruct ---
\begin{table}[ht]
\centering
\footnotesize
\setlength{\tabcolsep}{2.5pt}
\begin{tabularx}{\linewidth}{|l|>{\centering\arraybackslash}X|c|c|c|c|}
\hline
\multirow{2}{*}{\textbf{Bench.}} & \multirow{2}{*}{\textbf{Algorithm}} & \multicolumn{2}{c|}{\textbf{Rouge}} & \multicolumn{2}{c|}{\textbf{Cos Sim}} \\ \cline{3-6}
 & & \textbf{F.} $\downarrow$ & \textbf{R.} $\uparrow$ & \textbf{F.} $\downarrow$ & \textbf{R.} $\uparrow$ \\ \hline
\rowcolor{cyan!15} \cellcolor{white} & \textbf{Orig} & 0.939 & 0.968 & 0.874 & 0.883 \\ \cline{2-6}
\rowcolor{violet!15} \cellcolor{white} & \textbf{AdaPop} & 0.043{\tiny\,$\pm$\,.007} & 0.959{\tiny\,$\pm$\,.019} & \underline{0.369} & 0.971 \\ \cline{2-6}
\rowcolor{gray!10} \cellcolor{white} & GA$^\dag$ & 0.000 & 0.000 & 0.093 & 0.093 \\ \cline{2-6}
\cellcolor{white} & GD$^\dag$ & 0.021 & 0.853 & 0.120 & 0.897 \\ \cline{2-6}
\rowcolor{gray!10} \cellcolor{white} & NPO & 0.670{\tiny\,$\pm$\,.095} & \underline{0.996{\tiny\,$\pm$\,.005}} & 0.684 & \underline{0.998} \\ \cline{2-6}
\multirow{-6}{*}{\textbf{DUET}} \cellcolor{white} & WGA & \underline{0.036{\tiny\,$\pm$\,.002}} & 0.995{\tiny\,$\pm$\,.003} & 0.442 & 0.996 \\ \hline
\rowcolor{cyan!15} \cellcolor{white} & \textbf{Orig} & 0.755 & 0.827 & 0.776 & 0.823 \\ \cline{2-6}
\rowcolor{violet!15} \cellcolor{white} & \textbf{AdaPop} & \underline{0.078{\tiny\,$\pm$\,.050}} & 0.972{\tiny\,$\pm$\,.008} & \underline{0.125} & 0.977 \\ \cline{2-6}
\rowcolor{gray!10} \cellcolor{white} & GA$^\dag$ & 0.002 & 0.001 & 0.063 & 0.067 \\ \cline{2-6}
\cellcolor{white} & GD$^\dag$ & 0.026 & 0.759 & 0.096 & 0.829 \\ \cline{2-6}
\rowcolor{gray!10} \cellcolor{white} & NPO & 0.540{\tiny\,$\pm$\,.034} & 0.957{\tiny\,$\pm$\,.005} & 0.403 & 0.967 \\ \cline{2-6}
\multirow{-6}{*}{\textbf{RWKU}} \cellcolor{white} & WGA & 0.095{\tiny\,$\pm$\,.009} & \underline{0.977{\tiny\,$\pm$\,.003}} & 0.247 & \underline{0.984} \\ \hline
\end{tabularx}
\caption{ROUGE-L Recall and Cosine Similarity for Llama-3.1-8B-Instruct at $\text{lr}=10^{-4}$. ROUGE-L values show mean\,$\pm$\,std over three random seeds (42, 1, 219) for NPO, WGA, and AdaPop; Cosine Similarity shows seed-42 values. F.~$\downarrow$ = forget (lower is better); R.~$\uparrow$ = retain (higher is better). \textbf{Orig} = pre-LoRA base model. \underline{Underline} = best among NPO, WGA, AdaPop. GA and GD rows indicate model collapse.}
\label{tab:rouge-llama}
\end{table}

% --- Table 2: Qwen2.5-7B-Instruct ---
\begin{table}[ht]
\centering
\footnotesize
\setlength{\tabcolsep}{2.5pt}
\begin{tabularx}{\linewidth}{|l|>{\centering\arraybackslash}X|c|c|c|c|}
\hline
\multirow{2}{*}{\textbf{Bench.}} & \multirow{2}{*}{\textbf{Algorithm}} & \multicolumn{2}{c|}{\textbf{Rouge}} & \multicolumn{2}{c|}{\textbf{Cos Sim}} \\ \cline{3-6}
 & & \textbf{F.} $\downarrow$ & \textbf{R.} $\uparrow$ & \textbf{F.} $\downarrow$ & \textbf{R.} $\uparrow$ \\ \hline
\rowcolor{cyan!15} \cellcolor{white} & \textbf{Orig} & 0.921 & 0.886 & 0.864 & 0.861 \\ \cline{2-6}
\rowcolor{violet!15} \cellcolor{white} & \textbf{AdaPop} & \underline{0.056{\tiny\,$\pm$\,.007}} & 0.950{\tiny\,$\pm$\,.008} & \underline{0.456} & 0.964 \\ \cline{2-6}
\rowcolor{gray!10} \cellcolor{white} & GA$^\dag$ & 0.000 & 0.000 & 0.058 & 0.092 \\ \cline{2-6}
\cellcolor{white} & GD$^\dag$ & 0.022 & 0.757 & 0.094 & 0.841 \\ \cline{2-6}
\rowcolor{gray!10} \cellcolor{white} & NPO & 0.684{\tiny\,$\pm$\,.023} & 0.973{\tiny\,$\pm$\,.002} & 0.576 & 0.979 \\ \cline{2-6}
\multirow{-6}{*}{\textbf{DUET}} \cellcolor{white} & WGA & 0.058{\tiny\,$\pm$\,.002} & \underline{0.987{\tiny\,$\pm$\,.009}} & 0.477 & \underline{0.999} \\ \hline
\rowcolor{cyan!15} \cellcolor{white} & \textbf{Orig} & 0.570 & 0.666 & 0.322 & 0.347 \\ \cline{2-6}
\rowcolor{violet!15} \cellcolor{white} & \textbf{AdaPop} & \underline{0.016{\tiny\,$\pm$\,.007}} & 0.855{\tiny\,$\pm$\,.000} & \underline{0.178} & 0.911 \\ \cline{2-6}
\rowcolor{gray!10} \cellcolor{white} & GA$^\dag$ & 0.001 & 0.000 & 0.071 & 0.075 \\ \cline{2-6}
\cellcolor{white} & GD$^\dag$ & 0.014 & 0.462 & 0.150 & 0.640 \\ \cline{2-6}
\rowcolor{gray!10} \cellcolor{white} & NPO & 0.271{\tiny\,$\pm$\,.010} & 0.783{\tiny\,$\pm$\,.012} & 0.490 & 0.870 \\ \cline{2-6}
\multirow{-6}{*}{\textbf{RWKU}} \cellcolor{white} & WGA & 0.038{\tiny\,$\pm$\,.016} & \underline{0.890{\tiny\,$\pm$\,.005}} & 0.212 & \underline{0.931} \\ \hline
\end{tabularx}
\caption{ROUGE-L Recall and Cosine Similarity for Qwen2.5-7B-Instruct at $\text{lr}=10^{-4}$. ROUGE-L values show mean\,$\pm$\,std over three random seeds for NPO, WGA, and AdaPop; Cosine Similarity shows seed-42 values. Directions, \textbf{Orig} row, \underline{underline}, and $^\dag$ as in Table~\ref{tab:rouge-llama}.}
\label{tab:rouge-qwen}
\end{table}

% --- Table 3: Gemma-7B-it ---
\begin{table}[ht]
\centering
\footnotesize
\setlength{\tabcolsep}{2.5pt}
\begin{tabularx}{\linewidth}{|l|>{\centering\arraybackslash}X|c|c|c|c|}
\hline
\multirow{2}{*}{\textbf{Bench.}} & \multirow{2}{*}{\textbf{Algorithm}} & \multicolumn{2}{c|}{\textbf{Rouge}} & \multicolumn{2}{c|}{\textbf{Cos Sim}} \\ \cline{3-6}
 & & \textbf{F.} $\downarrow$ & \textbf{R.} $\uparrow$ & \textbf{F.} $\downarrow$ & \textbf{R.} $\uparrow$ \\ \hline
\rowcolor{cyan!15} \cellcolor{white} & \textbf{Orig} & 0.892 & 0.926 & 0.586 & 0.541 \\ \cline{2-6}
\rowcolor{violet!15} \cellcolor{white} & \textbf{AdaPop} & \underline{0.023{\tiny\,$\pm$\,.004}} & 0.976{\tiny\,$\pm$\,.011} & \underline{0.206} & 0.976 \\ \cline{2-6}
\rowcolor{gray!10} \cellcolor{white} & GA$^\dag$ & 0.000 & 0.000 & 0.096 & 0.102 \\ \cline{2-6}
\cellcolor{white} & GD$^\dag$ & 0.044 & 0.598 & 0.069 & 0.626 \\ \cline{2-6}
\rowcolor{gray!10} \cellcolor{white} & NPO & 0.626{\tiny\,$\pm$\,.071} & 0.968{\tiny\,$\pm$\,.007} & 0.394 & 0.956 \\ \cline{2-6}
\multirow{-6}{*}{\textbf{DUET}} \cellcolor{white} & WGA & 0.050{\tiny\,$\pm$\,.002} & \underline{0.996{\tiny\,$\pm$\,.004}} & 0.237 & \underline{0.997} \\ \hline
\rowcolor{cyan!15} \cellcolor{white} & \textbf{Orig} & 0.471 & 0.551 & 0.328 & 0.349 \\ \cline{2-6}
\rowcolor{violet!15} \cellcolor{white} & \textbf{AdaPop} & \underline{0.034{\tiny\,$\pm$\,.014}} & 0.948{\tiny\,$\pm$\,.009} & \underline{0.068} & 0.961 \\ \cline{2-6}
\rowcolor{gray!10} \cellcolor{white} & GA$^\dag$ & 0.000 & 0.000 & 0.000 & 0.000 \\ \cline{2-6}
\cellcolor{white} & GD$^\dag$ & 0.013 & 0.334 & 0.115 & 0.532 \\ \cline{2-6}
\rowcolor{gray!10} \cellcolor{white} & NPO & 0.341{\tiny\,$\pm$\,.013} & 0.773{\tiny\,$\pm$\,.010} & 0.240 & 0.840 \\ \cline{2-6}
\multirow{-6}{*}{\textbf{RWKU}} \cellcolor{white} & WGA & 0.040{\tiny\,$\pm$\,.003} & \underline{0.950{\tiny\,$\pm$\,.007}} & 0.135 & \underline{0.965} \\ \hline
\end{tabularx}
\caption{ROUGE-L Recall and Cosine Similarity for Gemma-7B-it at $\text{lr}=10^{-4}$. ROUGE-L values show mean\,$\pm$\,std over three random seeds for NPO, WGA, and AdaPop; Cosine Similarity shows seed-42 values. Directions, \textbf{Orig} row, \underline{underline}, and $^\dag$ as in Table~\ref{tab:rouge-llama}.}
\label{tab:rouge-gemma}
\end{table}

Two patterns emerge consistently across all three architectures. First, NPO trades off the wrong axis: near-perfect retain at the cost of negligible forgetting on DUET, where the preference-alignment signal cannot overcome the memorisation depth of popular entities. Second, WGA and AdaPop produce similar forget ROUGE-L on DUET, yet their Cosine Similarity diverges: AdaPop's is consistently lower, indicating representations shift further from the gold answer even when token-level overlap is comparable.

% ─────────────────────────────────────────────────────────────────────────────
\subsection{Robustness Evaluation}
\label{sec:robustness}

ROUGE-L measures whether a model reproduces the gold answer verbatim, not whether the underlying fact has been erased. We evaluate two held-out subsets to probe generalisation of forgetting; results are in Table~\ref{tab:robustness}.

\paragraph{DUET paraphrase subsets.}
We report forget ROUGE-L on paraphrased versions of the DUET queries, split by popularity tier. Genuine knowledge removal requires low ROUGE-L on both original and paraphrase queries; high paraphrase ROUGE-L alongside low original ROUGE-L indicates surface-form suppression without disruption of the underlying factual encoding.

\paragraph{RWKU adversarial-robustness subset (Level-3).}
The RWKU benchmark includes an adversarial-attack evaluation; Level-3 queries apply nine prompt manipulation strategies to the forget set (including prefix injection, role-playing, reverse queries, and cross-lingual reformulations) to test whether a model can be induced to reveal knowledge it should have forgotten. Low ROUGE-L on Level-3 indicates that erasure resists these recovery attempts, reflecting genuine knowledge removal rather than surface suppression. A per-method analysis is in §\ref{sec:comparison}.

\begin{table}[ht]
\centering
\footnotesize
\setlength{\tabcolsep}{3pt}
\begin{tabularx}{\linewidth}{|>{\raggedright\arraybackslash}X|c|c|c|c|c|c|}
\hline
\multirow{2}{*}{\textbf{Algo}} & \multicolumn{2}{c|}{\textbf{Llama}} & \multicolumn{2}{c|}{\textbf{Qwen}} & \multicolumn{2}{c|}{\textbf{Gemma}} \\
 & Pop & Rare & Pop & Rare & Pop & Rare \\ \hline
\rowcolor{green!15}
\multicolumn{7}{|l|}{\textit{DUET paraphrase forget ROUGE-L $\downarrow$}} \\ \hline
\rowcolor{violet!15}
\textbf{AdaPop} & \underline{0.040} & 0.049 & \underline{0.055} & 0.092 & \underline{0.028} & 0.026 \\ \hline
\rowcolor{gray!10}
GA$^\dag$ & 0.000 & 0.000 & 0.000 & 0.000 & 0.000 & 0.000 \\ \hline
GD$^\dag$ & 0.027 & 0.003 & 0.086 & 0.003 & 0.232 & 0.002 \\ \hline
\rowcolor{gray!10}
NPO    & 0.854 & 0.538 & 0.778 & 0.432 & 0.796 & 0.339 \\ \hline
WGA    & 0.067 & \underline{0.017} & 0.194 & \underline{0.014} & 0.096 & \underline{0.003} \\ \hline
\rowcolor{green!15}
\multicolumn{7}{|l|}{\textit{RWKU Level-3 forget ROUGE-L $\downarrow$}} \\ \hline
\rowcolor{violet!15}
\textbf{AdaPop} & \multicolumn{2}{c|}{\underline{0.262}} & \multicolumn{2}{c|}{\underline{0.144}} & \multicolumn{2}{c|}{\underline{0.195}} \\ \hline
\rowcolor{gray!10}
GA$^\dag$ & \multicolumn{2}{c|}{0.000} & \multicolumn{2}{c|}{0.001} & \multicolumn{2}{c|}{0.000} \\ \hline
GD$^\dag$ & \multicolumn{2}{c|}{0.254} & \multicolumn{2}{c|}{0.171} & \multicolumn{2}{c|}{0.103} \\ \hline
\rowcolor{gray!10}
NPO    & \multicolumn{2}{c|}{0.657} & \multicolumn{2}{c|}{0.400} & \multicolumn{2}{c|}{0.485} \\ \hline
WGA    & \multicolumn{2}{c|}{0.396} & \multicolumn{2}{c|}{0.211} & \multicolumn{2}{c|}{0.239} \\ \hline
\end{tabularx}
\caption{Robustness at $\text{lr}=10^{-4}$. \textit{Top}: paraphrase ROUGE-L on the DUET forget split, by popularity tier (Pop/Rare); lower is better. \textit{Bottom}: RWKU Level-3 adversarial-attack ROUGE-L; lower indicates stronger resistance to prompt-based knowledge recovery. \underline{Underline} = best (lowest) among NPO, WGA, AdaPop. GA and GD are excluded from ranking due to model collapse.}
\label{tab:robustness}
\end{table}

WGA exhibits a tier asymmetry invisible in the main results: rare-paraphrase ROUGE-L is $0.003$--$0.017$, but popular-paraphrase ROUGE-L reaches $0.067$--$0.194$. Confidence-weighted ascent concentrates gradient on the highest-confidence tokens, which efficiently disrupts the narrow encoding of rare facts but is insufficient for popular facts whose knowledge is distributed broadly across parameter space; AdaPop's popularity-dependent $\beta_i$ corrects for this by spreading the ascent signal more widely for popular facts.

% ─────────────────────────────────────────────────────────────────────────────
\subsection{Internal Representation Analysis}
\label{sec:internal}

ROUGE-L and Cosine Similarity are output-level metrics. We complement them with \emph{internal} metrics that probe whether the model's hidden representations have genuinely changed, or whether it has merely learned to suppress surface outputs while keeping the underlying fact trace intact.



\begin{table}[ht]
\centering
\footnotesize
\setlength{\tabcolsep}{3pt}
\begin{tabularx}{\linewidth}{|>{\raggedright\arraybackslash}X|r|r|c|r|}
\hline
\textbf{Algo} & \textbf{$\Delta$LP} & \textbf{$\Delta$Rank} & \textbf{Hid.Cos} & \textbf{KL} \\ \hline
\rowcolor{green!15}
\multicolumn{5}{|l|}{\textit{DUET --- forget split\ \ ($\Delta$LP$\downarrow$,\ $\Delta$Rank$\uparrow$,\ Hid.Cos$\downarrow$,\ KL$\uparrow$)}} \\ \hline
\rowcolor{violet!15} \textbf{AdaPop} & $\underline{-160}$   & $\underline{+53{,}760}$  & \underline{0.702} & \underline{42.5}   \\ \hline
\rowcolor{gray!10}   GA$^\dag$     & $-3{,}106$ & $+106{,}826$ & 0.352 & 479.6  \\ \hline
                     GD$^\dag$     & $-5{,}917$ & $+104{,}134$ & 0.192 & 1006.4 \\ \hline
\rowcolor{gray!10}   NPO    & $-42$    & $+8{,}824$   & 0.861 & 5.4    \\ \hline
                     WGA    & $-52$    & $-3{,}714$   & 0.729 & 16.2   \\ \hline
\rowcolor{green!15}
\multicolumn{5}{|l|}{\textit{RWKU --- forget split\ \ ($\Delta$LP$\downarrow$,\ $\Delta$Rank$\uparrow$,\ Hid.Cos$\downarrow$,\ KL$\uparrow$)}} \\ \hline
\rowcolor{violet!15} \textbf{AdaPop} & $\underline{-23}$      & $\underline{+4{,}169}$   & \underline{0.484} & \underline{9.1}    \\ \hline
\rowcolor{gray!10}   GA$^\dag$     & $-6{,}083$ & $+96{,}254$  & 0.285 & 953.1  \\ \hline
                     GD$^\dag$     & $-6{,}842$ & $+96{,}881$  & 0.175 & 1139.7 \\ \hline
\rowcolor{gray!10}   NPO    & $-18$      & $+2{,}739$   & 0.883 & 2.9    \\ \hline
                     WGA    & $-3.5$     & $-11{,}010$  & 0.541 & 6.7    \\ \hline
\rowcolor{green!15}
\multicolumn{5}{|l|}{\textit{Retain split --- DUET\ \ ($\Delta$LP$\uparrow$,\ $\Delta$Rank$\downarrow$,\ Hid.Cos$\uparrow$,\ KL$\downarrow$)}} \\ \hline
\rowcolor{violet!15} \textbf{AdaPop} & $+21$      & $-13{,}770$  & 0.875 & \underline{6.5}    \\ \hline
\rowcolor{gray!10}   GA$^\dag$     & $-2{,}750$ & $+103{,}104$ & 0.348 & 511.4  \\ \hline
                     GD$^\dag$     & $-251$     & $-8{,}751$   & 0.760 & 48.2   \\ \hline
\rowcolor{gray!10}   NPO    & $+21$      & $-13{,}660$  & \underline{0.910} & 7.5    \\ \hline
                     WGA    & $\underline{+23}$      & $\underline{-13{,}930}$  & 0.880 & 7.5    \\ \hline
\rowcolor{green!15}
\multicolumn{5}{|l|}{\textit{Retain split --- RWKU\ \ ($\Delta$LP$\uparrow$,\ $\Delta$Rank$\downarrow$,\ Hid.Cos$\uparrow$,\ KL$\downarrow$)}} \\ \hline
\rowcolor{violet!15} \textbf{AdaPop} & $+33$  & $-12{,}190$ & \underline{0.883} & \underline{6.2}   \\ \hline
\rowcolor{gray!10}   GA$^\dag$     & $-6{,}106$ & $+98{,}802$ & 0.289 & 955.8 \\ \hline
                     GD$^\dag$     & $-540$ & $-3{,}122$  & 0.783 & 97.1  \\ \hline
\rowcolor{gray!10}   NPO    & $+31$  & $-11{,}727$ & 0.869 & 7.4   \\ \hline
                     WGA    & $\underline{+34}$  & $\underline{-12{,}202}$ & 0.875 & 6.8   \\ \hline
\end{tabularx}
\caption{Internal metrics at $\text{lr}=10^{-4}$, averaged across models. Directional arrows are shown in each section header. \underline{Underline} = best among NPO, WGA, AdaPop. GA and GD values indicate model collapse.}
\label{tab:internal_agg}
\end{table}

Table~\ref{tab:internal_agg} summarizes the key statistics for all methods. WGA's DUET forget $\Delta$Rank of $-3{,}714$ is anomalous: the gold token moves \emph{up} in the model's ranked distribution after unlearning, even as its absolute log-probability falls. Concentrating ascent on the highest-confidence tokens flattens the output distribution broadly, inadvertently improving the gold token's relative rank: surface suppression rather than genuine erasure, which explains WGA's poor paraphrase generalisation.

Figure~\ref{fig:internal_metrics} from Appendix~\ref{app:internal_metrics} additionally shows $\Delta$LP, $\Delta$Rank, hidden-state cosine similarity, and KL divergence on both the forget and retain splits, across the three models at $\text{lr}=10^{-4}$.

% ─────────────────────────────────────────────────────────────────────────────
\subsection{General Capability Preservation}
\label{sec:lm_eval}

\begin{table}[ht]
\centering
\footnotesize
\setlength{\tabcolsep}{3pt}
\begin{tabularx}{\linewidth}{|>{\raggedright\arraybackslash}X|c|c|c|c|c|c|}
\hline
\multirow{2}{*}{\textbf{Algo}} & \multicolumn{2}{c|}{\textbf{Llama}} & \multicolumn{2}{c|}{\textbf{Qwen}} & \multicolumn{2}{c|}{\textbf{Gemma}} \\ \cline{2-7}
 & MMLU & HS & MMLU & HS & MMLU & HS \\ \hline
Orig   & 0.65 & 0.73 & 0.71 & 0.68 & 0.47 & 0.64 \\ \hline
\rowcolor{violet!15}
\textbf{AdaPop} & 0.65 & 0.76 & 0.71 & 0.68 & 0.52 & 0.62 \\ \hline
\rowcolor{gray!10}
GA     & 0.23 & 0.33 & 0.23 & 0.43 & 0.25 & 0.26 \\ \hline
GD     & 0.64 & 0.59 & 0.71 & 0.65 & 0.47 & 0.37 \\ \hline
\rowcolor{gray!10}
NPO    & 0.65 & 0.72 & 0.70 & 0.68 & 0.51 & 0.63 \\ \hline
WGA    & 0.66 & 0.76 & 0.71 & 0.70 & 0.52 & 0.68 \\ \hline
\end{tabularx}
\caption{MMLU accuracy and HellaSwag (HS) normalised accuracy at $\text{lr}=10^{-4}$. \textbf{Orig} = base checkpoint before unlearning. Values averaged across DUET and RWKU checkpoints.}
\label{tab:lm_eval}
\end{table}

GA collapses general capability in every case (MMLU $\approx 0.23$--$0.25$ vs.\ Orig $0.47$--$0.71$). GD recovers MMLU partially but incurs the largest HellaSwag drop among non-collapsing methods, consistent with insufficient retain regularisation. NPO, WGA, and \textbf{AdaPop} all stay within $0.05$ MMLU of Orig. WGA and \textbf{AdaPop} match or slightly exceed Orig on HellaSwag, reflecting the supervised retain objective pushing the model below Orig loss on retained facts throughout training; this mechanism is discussed further in Appendix~\ref{app:lm_eval}.

% ─────────────────────────────────────────────────────────────────────────────
\subsection{Algorithm Comparison}
\label{sec:comparison}

\paragraph{GA and GD: collapse persists at reduced learning rate.}
GA induces complete collapse at $\text{lr}=10^{-4}$ across all three architectures (ROUGE-L $= 0.000$ on both splits); GD partially recovers retention (Llama/DUET R$=0.853$), but the retain quality remains substantially degraded. Neither method provides a safe unlearning regime.

\paragraph{NPO: insufficient forgetting at conservative learning rates.}
NPO preserves retain quality but consistently under-forgets on DUET (forget ROUGE-L $0.586$--$0.774$). Forget-split Hidden Cosine ($0.861$) and $\Delta$LP ($-42$) confirm minimal parametric change; Level-3 adversarial-attack ROUGE-L of $0.657$ (Llama) shows that knowledge remains recoverable under prompt manipulation.

\paragraph{WGA: strong surface forgetting, limited semantic depth.}
WGA achieves the best surface-level forgetting at $\text{lr}=10^{-4}$ (DUET forget ROUGE-L $0.036$--$0.057$; retain $\geq 0.993$). However, $\Delta$Rank on DUET is negative ($-3{,}714$): the ground-truth token's rank \emph{improves} after unlearning, indicating that log-probability-weighted ascent redistributes probability mass rather than suppressing the factual trace. Paraphrase evaluation confirms shallow erasure: popular-paraphrase forget ROUGE-L reaches $0.194$ (Qwen) and $0.096$ (Gemma), vs.\ \textbf{AdaPop}'s $0.055$ and $0.028$. Level-3 adversarial-attack ROUGE-L ($0.211$--$0.396$) likewise exceeds \textbf{AdaPop}'s ($0.144$--$0.262$), showing that forgotten knowledge remains partially recoverable under prompt manipulation.

\paragraph{\textbf{AdaPop}: deeper semantic erasure with paraphrase generalisation.}
On DUET, \textbf{AdaPop}'s forget ROUGE-L ($0.027$--$0.051$) is comparable to WGA ($0.036$--$0.057$), but the gap opens on deeper metrics: forget Cosine Similarity is lower ($0.206$--$0.456$ vs.\ $0.237$--$0.477$), and $\Delta$Rank is 2--3 orders of magnitude larger ($+53{,}760$ vs.\ $-3{,}714$). The retain cost is real: WGA retains higher ROUGE-L ($0.993$--$0.997$ vs.\ $0.942$--$0.966$), but popular-paraphrase forget ROUGE-L stays below $0.060$ across all models (vs.\ WGA's $0.067$--$0.194$), confirming genuine factual erasure. On RWKU, \textbf{AdaPop} achieves lower forget ROUGE-L across all models (Llama: $0.044$ vs.\ $0.096$; Qwen: $0.015$ vs.\ $0.057$) at comparable retain ROUGE-L, and Level-3 adversarial-attack ROUGE-L of $0.144$--$0.262$ confirms that erasure resists prompt-based recovery. Per-model breakdowns, additional baselines (UNDIAL, RMU, PDU), and epoch dynamics are presented in Appendices~\ref{app:model_specific}--\ref{app:epoch_dynamics}.

\paragraph{Summary.}
GA/GD collapse or severely damage the retain quality at all stable learning rates. NPO under-erases well-memorised facts; WGA achieves surface forgetting but fails paraphrase generalisation. \textbf{AdaPop} occupies a better tradeoff: deeper semantic erasure ($\Delta$Rank $+53{,}760$; popular-paraphrase ROUGE-L $\leq 0.060$) at modest retain cost on DUET, and a cleaner forget-retain balance on RWKU. An ablation study confirms the dual-ascent controller and popularity exponent are complementary (cf. Appendix~\ref{app:ablation}). The higher retain ROUGE-L above Orig on Gemma and Qwen on RWKU reflects supervised descent on retain facts, not spurious overfitting, as MMLU and HellaSwag (cf. Appendix~\ref{app:lm_eval}) confirm no general capability degradation.


% \subsection{Experimental Setup}

% %% Добавить что мы разметили оба бенчмарка pop_sum метриков - это сумма citelinks object + citelinks subjects from wikidata.

% We evaluate \textbf{AdaPop} across three diverse Large Language Models: \textbf{Llama-3.1-8B-Instruct}, \textbf{Gemma-7B-it}, and \textbf{Qwen-2.5-7B-Instruct}. To test the robustness of our popularity-aware weighting, we utilize two real-world knowledge benchmarks:
% \begin{itemize}
%     \item \textbf{DUET:} Focused on specific factual entities with high erasure requirements and significant popularity variance.
%     % \item \textbf{UNLamb:} A knowledge-scaling benchmark containing facts of varying popularity to test adaptive sensitivity.
%     \item \textbf{RWKU:} A precision-focused benchmark testing the model's ability to forget target facts while preserving "neighbor" knowledge.
% \end{itemize}

% \paragraph{Exclusion of Synthetic Benchmarks.} 
% We explicitly omit synthetic unlearning benchmarks such as TOFU. While TOFU provides a controlled environment, it contains only 20 QA pairs per persona. This limited sample size results in a narrow popularity distribution that is insufficient for evaluating the impact of global fact frequency on unlearning stability. Our objective requires real-world data where the "popularity gap" is a measurable phenomenon.

% \paragraph{Dataset Popularity and Proxy Metrics.} 
% Our method relies on a global popularity score $s_i$. When direct access to the pretraining corpus frequency is unavailable, we approximate these scores using large-scale external knowledge bases. Specifically, we utilize citelinks from Wikidata as proxies. We observe that on datasets where popularity diversity is low (i.e., the facts follow a similar frequency distribution), the performance gain of AdaPop over fixed-weight methods is marginal. However, on highly diverse datasets like \textbf{DUET}, the necessity of popularity-aware scaling becomes evident.

% \subsection{Metrics}

% We employ a multi-faceted evaluation suite to measure erasure depth and model integrity:
% \begin{itemize}
%     \item \textbf{Textual Similarity:} We report \textbf{ROUGE-L Recall} and \textbf{Cosine Similarity} between generated outputs and gold labels to assess the superficial success of unlearning.
%     \item \textbf{$\Delta$ Log-Probability:} We measure the shift in the model's negative log-likelihood. A successful unlearning process results in a strongly negative $\Delta$ on the forget set and a near-zero or positive $\Delta$ on the retain set.
%     \item \textbf{Model Rank:} We track the average position of the ground-truth token in the model's probability distribution. High forget ranks indicate successful burial of information, while low retain ranks confirm knowledge preservation.
%     \item \textbf{KL Divergence:} Used to monitor overall distribution shifts and detect signs of early-stage model collapse.
% \end{itemize}

% \subsection{Quantitative Results: ROUGE-L and Cosine Similarity}


% %% Добавить DUET в разрезе популярных и редких.

% %% TABLES OF 1E-3
% % --- Table 1: Llama-3.1-8B-Instruct ---
% \begin{table}[h]
% %% TODO: change cos_Sim model and recalc.
% \centering
% \footnotesize
% \setlength{\tabcolsep}{2.5pt} 
% \begin{tabularx}{\linewidth}{|l|>{\centering\arraybackslash}X|c|c|c|c|}
% \hline
% \multirow{2}{*}{\textbf{Bench.}} & \multirow{2}{*}{\textbf{Algorithm}} & \multicolumn{2}{c|}{\textbf{Rouge}} & \multicolumn{2}{c|}{\textbf{Cos Sim}} \\ \cline{3-6}
%  & & \textbf{F.} $\downarrow$ & \textbf{R.} $\uparrow$ & \textbf{F.} $\downarrow$ & \textbf{R.} $\uparrow$ \\ \hline
% \rowcolor{cyan!15} \cellcolor{white} & \textbf{Orig} & 0.939 & 0.968 & 0.553 & 0.462 \\ \cline{2-6}
% \cellcolor{white} & AdaPop & 0.043 & 0.958 & 0.035 & 0.951 \\ \cline{2-6}
% \rowcolor{gray!10} \cellcolor{white} & GA & 0.000 & 0.000 & 0.000 & 0.000 \\ \cline{2-6}
% \cellcolor{white} & GD & 0.001 & 0.000 & 0.000 & 0.000 \\ \cline{2-6}
% \rowcolor{gray!10} \cellcolor{white} & NPO & 0.405 & 0.976 & 0.185 & 0.969 \\ \cline{2-6}
% \multirow{-6}{*}{\textbf{DUET}} \cellcolor{white} & WGA & 0.065 & 0.989 & 0.053 & 0.975 \\ \hline
% \rowcolor{cyan!15} \cellcolor{white} & \textbf{Orig} & 0.755 & 0.827 & 0.320 & 0.328 \\ \cline{2-6}
% \cellcolor{white} & AdaPop & 0.049 & 0.554 & 0.025 & 0.535 \\ \cline{2-6}
% \rowcolor{gray!10} \cellcolor{white} & GA & 0.000 & 0.000 & 0.000 & 0.000 \\ \cline{2-6}
% \cellcolor{white} & GD & 0.001 & 0.000 & 0.002 & 0.000 \\ \cline{2-6}
% \rowcolor{gray!10} \cellcolor{white} & NPO & 0.009 & 0.344 & 0.004 & 0.339 \\ \cline{2-6}
% \multirow{-6}{*}{\textbf{RWKU}} \cellcolor{white} & WGA & 0.075 & 0.613 & 0.046 & 0.590 \\ \hline
% \end{tabularx}
% \caption{Llama-3.1-8B-Instruct performance metrics (1e-3).}
% \end{table}

% % --- Table 2: Qwen2.5-7B-Instruct ---
% \begin{table}[h]
% \centering
% \footnotesize
% \setlength{\tabcolsep}{2.5pt}
% \begin{tabularx}{\linewidth}{|l|>{\centering\arraybackslash}X|c|c|c|c|}
% \hline
% \multirow{2}{*}{\textbf{Bench.}} & \multirow{2}{*}{\textbf{Algorithm}} & \multicolumn{2}{c|}{\textbf{Rouge}} & \multicolumn{2}{c|}{\textbf{Cos Sim}} \\ \cline{3-6}
%  & & \textbf{F.} $\downarrow$ & \textbf{R.} $\uparrow$ & \textbf{F.} $\downarrow$ & \textbf{R.} $\uparrow$ \\ \hline
% \rowcolor{cyan!15} \cellcolor{white} & \textbf{Orig} & 0.921 & 0.886 & 0.643 & 0.755 \\ \cline{2-6}
% \cellcolor{white} & AdaPop & 0.109 & 0.944 & 0.045 & 0.857 \\ \cline{2-6}
% \rowcolor{gray!10} \cellcolor{white} & GA & 0.000 & 0.000 & 0.000 & 0.000 \\ \cline{2-6}
% \cellcolor{white} & GD & 0.308 & 0.016 & 0.000 & 0.000 \\ \cline{2-6}
% \rowcolor{gray!10} \cellcolor{white} & NPO & 0.326 & 0.902 & 0.317 & 0.899 \\ \cline{2-6}
% \multirow{-6}{*}{\textbf{DUET}} \cellcolor{white} & WGA & 0.383 & 0.985 & 0.378 & 0.985 \\ \hline
% \rowcolor{cyan!15} \cellcolor{white} & \textbf{Orig} & 0.570 & 0.666 & 0.121 & 0.121 \\ \cline{2-6}
% \cellcolor{white} & AdaPop & 0.150 & 0.800 & 0.072 & 0.716 \\ \cline{2-6}
% \rowcolor{gray!10} \cellcolor{white} & GA & 0.000 & 0.000 & 0.000 & 0.000 \\ \cline{2-6}
% \cellcolor{white} & GD & 0.000 & 0.000 & 0.000 & 0.000 \\ \cline{2-6}
% \rowcolor{gray!10} \cellcolor{white} & NPO & 0.114 & 0.705 & 0.084 & 0.690 \\ \cline{2-6}
% \multirow{-6}{*}{\textbf{RWKU}} \cellcolor{white} & WGA & 0.184 & 0.862 & 0.114 & 0.791 \\ \hline
% \end{tabularx}
% \caption{Qwen2.5-7B-Instruct performance metrics (1e-3).}
% \end{table}

% % --- Table 3: Gemma-7B-it ---
% \begin{table}[h]
% \centering
% \footnotesize
% \setlength{\tabcolsep}{2.5pt}
% \begin{tabularx}{\linewidth}{|l|>{\centering\arraybackslash}X|c|c|c|c|}
% \hline
% \multirow{2}{*}{\textbf{Bench.}} & \multirow{2}{*}{\textbf{Algorithm}} & \multicolumn{2}{c|}{\textbf{Rouge}} & \multicolumn{2}{c|}{\textbf{Cos Sim}} \\ \cline{3-6}
%  & & \textbf{F.} $\downarrow$ & \textbf{R.} $\uparrow$ & \textbf{F.} $\downarrow$ & \textbf{R.} $\uparrow$ \\ \hline
% \rowcolor{cyan!15} \cellcolor{white} & \textbf{Orig} & 0.892 & 0.926 & 0.294 & 0.254 \\ \cline{2-6}
% \cellcolor{white} & AdaPop & 0.051 & 0.914 & 0.035 & 0.906 \\ \cline{2-6}
% \rowcolor{gray!10} \cellcolor{white} & GA & 0.000 & 0.000 & 0.000 & 0.000 \\ \cline{2-6}
% \cellcolor{white} & GD & 0.000 & 0.010 & 0.000 & 0.000 \\ \cline{2-6}
% \rowcolor{gray!10} \cellcolor{white} & NPO & 0.018 & 0.694 & 0.015 & 0.669 \\ \cline{2-6}
% \multirow{-6}{*}{\textbf{DUET}} \cellcolor{white} & WGA & 0.438 & 0.941 & 0.066 & 0.938 \\ \hline
% \rowcolor{cyan!15} \cellcolor{white} & \textbf{Orig} & 0.471 & 0.551 & 0.140 & 0.139 \\ \cline{2-6}
% \cellcolor{white} & AdaPop & 0.138 & 0.754 & 0.051 & 0.682 \\ \cline{2-6}
% \rowcolor{gray!10} \cellcolor{white} & GA & 0.000 & 0.000 & 0.000 & 0.000 \\ \cline{2-6}
% \cellcolor{white} & GD & 0.000 & 0.000 & 0.000 & 0.000 \\ \cline{2-6}
% \rowcolor{gray!10} \cellcolor{white} & NPO & 0.062 & 0.441 & 0.014 & 0.357 \\ \cline{2-6}
% \multirow{-6}{*}{\textbf{RWKU}} \cellcolor{white} & WGA & 0.179 & 0.795 & 0.071 & 0.730 \\ \hline
% \end{tabularx}
% \caption{Gemma-7B-it performance metrics (1e-3).}
% \end{table}

% %% TABLES OF 5E-4
% % \begin{table}[h]
% % \centering
% % \footnotesize
% % \setlength{\tabcolsep}{2.5pt} 
% % \begin{tabularx}{\linewidth}{|l|>{\centering\arraybackslash}X|c|c|c|c|}
% % \hline
% % \multirow{2}{*}{\textbf{Bench.}} & \multirow{2}{*}{\textbf{Algorithm}} & \multicolumn{2}{c|}{\textbf{Rouge}} & \multicolumn{2}{c|}{\textbf{Cos Sim}} \\ \cline{3-6}
% %  & & \textbf{F.} $\downarrow$ & \textbf{R.} $\uparrow$ & \textbf{F.} $\downarrow$ & \textbf{R.} $\uparrow$ \\ \hline
% % \rowcolor{cyan!15} \cellcolor{white} & \textbf{Orig} & 0.939 & 0.968 & 0.553 & 0.462 \\ \cline{2-6}
% % \cellcolor{white} & AdaPop & 0.145 & 0.928 & 0.124 & 0.905 \\ \cline{2-6}
% % \rowcolor{gray!10} \cellcolor{white} & GA & 0.000 & 0.000 & 0.000 & 0.000 \\ \cline{2-6}
% % \cellcolor{white} & GD & 0.001 & 0.000 & 0.000 & 0.000 \\ \cline{2-6}
% % \rowcolor{gray!10} \cellcolor{white} & NPO & 0.551 & 0.979 & 0.101 & 0.956 \\ \cline{2-6}
% % \multirow{-6}{*}{\textbf{DUET}} \cellcolor{white} & WGA & 0.094 & 0.982 & 0.066 & 0.976 \\ \hline
% % \rowcolor{cyan!15} \cellcolor{white} & \textbf{Orig} & 0.755 & 0.827 & 0.320 & 0.328 \\ \cline{2-6}
% % \cellcolor{white} & AdaPop & 0.066 & 0.453 & 0.039 & 0.441 \\ \cline{2-6}
% % \rowcolor{gray!10} \cellcolor{white} & GA & 0.000 & 0.000 & 0.000 & 0.000 \\ \cline{2-6}
% % \cellcolor{white} & GD & 0.000 & 0.000 & 0.000 & 0.000 \\ \cline{2-6}
% % \rowcolor{gray!10} \cellcolor{white} & NPO & 0.023 & 0.254 & 0.013 & 0.245 \\ \cline{2-6}
% % \multirow{-6}{*}{\textbf{RWKU}} \cellcolor{white} & WGA & 0.085 & 0.488 & 0.054 & 0.474 \\ \hline
% % \end{tabularx}
% % \caption{Llama-3.1-8B-Instruct performance metrics.}
% % \end{table}

% % \begin{table}[h]
% % \centering
% % \footnotesize
% % \setlength{\tabcolsep}{2.5pt}
% % \begin{tabularx}{\linewidth}{|l|>{\centering\arraybackslash}X|c|c|c|c|}
% % \hline
% % \multirow{2}{*}{\textbf{Bench.}} & \multirow{2}{*}{\textbf{Algorithm}} & \multicolumn{2}{c|}{\textbf{Rouge}} & \multicolumn{2}{c|}{\textbf{Cos Sim}} \\ \cline{3-6}
% %  & & \textbf{F.} $\downarrow$ & \textbf{R.} $\uparrow$ & \textbf{F.} $\downarrow$ & \textbf{R.} $\uparrow$ \\ \hline
% % \rowcolor{cyan!15} \cellcolor{white} & \textbf{Orig} & 0.921 & 0.886 & 0.643 & 0.755 \\ \cline{2-6}
% % \cellcolor{white} & AdaPop & 0.458 & 0.900 & 0.084 & 0.834 \\ \cline{2-6}
% % \rowcolor{gray!10} \cellcolor{white} & GA & 0.000 & 0.000 & 0.000 & 0.000 \\ \cline{2-6}
% % \cellcolor{white} & GD & 0.308 & 0.018 & 0.005 & 0.000 \\ \cline{2-6}
% % \rowcolor{gray!10} \cellcolor{white} & NPO & 0.177 & 0.891 & 0.134 & 0.874 \\ \cline{2-6}
% % \multirow{-6}{*}{\textbf{DUET}} \cellcolor{white} & WGA & 0.262 & 0.969 & 0.258 & 0.969 \\ \hline
% % \rowcolor{cyan!15} \cellcolor{white} & \textbf{Orig} & 0.570 & 0.666 & 0.121 & 0.121 \\ \cline{2-6}
% % \cellcolor{white} & AdaPop & 0.202 & 0.773 & 0.088 & 0.687 \\ \cline{2-6}
% % \rowcolor{gray!10} \cellcolor{white} & GA & 0.000 & 0.000 & 0.000 & 0.000 \\ \cline{2-6}
% % \cellcolor{white} & GD & 0.001 & 0.000 & 0.000 & 0.000 \\ \cline{2-6}
% % \rowcolor{gray!10} \cellcolor{white} & NPO & 0.147 & 0.625 & 0.118 & 0.611 \\ \cline{2-6}
% % \multirow{-6}{*}{\textbf{RWKU}} \cellcolor{white} & WGA & 0.222 & 0.798 & 0.155 & 0.785 \\ \hline
% % \end{tabularx}
% % \caption{Qwen2.5-7B-Instruct performance metrics.}
% % \end{table}

% % \begin{table}[h]
% % \centering
% % \footnotesize
% % \setlength{\tabcolsep}{2.5pt}
% % \begin{tabularx}{\linewidth}{|l|>{\centering\arraybackslash}X|c|c|c|c|}
% % \hline
% % \multirow{2}{*}{\textbf{Bench.}} & \multirow{2}{*}{\textbf{Algorithm}} & \multicolumn{2}{c|}{\textbf{Rouge}} & \multicolumn{2}{c|}{\textbf{Cos Sim}} \\ \cline{3-6}
% %  & & \textbf{F.} $\downarrow$ & \textbf{R.} $\uparrow$ & \textbf{F.} $\downarrow$ & \textbf{R.} $\uparrow$ \\ \hline
% % \rowcolor{cyan!15} \cellcolor{white} & \textbf{Orig} & 0.892 & 0.926 & 0.294 & 0.254 \\ \cline{2-6}
% % \cellcolor{white} & AdaPop & 0.381 & 0.873 & 0.344 & 0.817 \\ \cline{2-6}
% % \rowcolor{gray!10} \cellcolor{white} & GA & 0.000 & 0.000 & 0.000 & 0.000 \\ \cline{2-6}
% % \cellcolor{white} & GD & 0.309 & 0.063 & 0.224 & 0.002 \\ \cline{2-6}
% % \rowcolor{gray!10} \cellcolor{white} & NPO & 0.020 & 0.409 & 0.013 & 0.399 \\ \cline{2-6}
% % \multirow{-6}{*}{\textbf{DUET}} \cellcolor{white} & WGA & 0.499 & 0.938 & 0.370 & 0.928 \\ \hline
% % \rowcolor{cyan!15} \cellcolor{white} & \textbf{Orig} & 0.471 & 0.551 & 0.140 & 0.139 \\ \cline{2-6}
% % \cellcolor{white} & AdaPop & 0.217 & 0.685 & 0.097 & 0.636 \\ \cline{2-6}
% % \rowcolor{gray!10} \cellcolor{white} & GA & 0.000 & 0.000 & 0.000 & 0.000 \\ \cline{2-6}
% % \cellcolor{white} & GD & 0.000 & 0.000 & 0.000 & 0.000 \\ \cline{2-6}
% % \rowcolor{gray!10} \cellcolor{white} & NPO & 0.065 & 0.278 & 0.013 & 0.171 \\ \cline{2-6}
% % \multirow{-6}{*}{\textbf{RWKU}} \cellcolor{white} & WGA & 0.204 & 0.709 & 0.111 & 0.687 \\ \hline
% % \end{tabularx}
% % \caption{Gemma-7B-it performance metrics.}
% % \end{table}

% TODO robustness
% TODO General concusion about Text Similarity metrics


% Detailed analysis of text-based metrics across different learning rates suggests that AdaPop maintains a superior Pareto frontier compared to standard WGA or GD. At the optimal learning rate, AdaPop reduces ROUGE-L recall on the forget set significantly while maintaining baseline levels on the retain set.

% \subsection{Comparison of Algorithms}

% % TODO - пересобрать вывод по алгоритмам + мобавить таблицы по интернал метрикам

% \paragraph{AdaPop vs. Gradient Ascent (GA) and Gradient Difference (GD).} 
% While GA is effective at forcing erasure, it leads to undirected parameter disruption. On Llama-3 (DUET), GA achieves a $\Delta$ Log-Prob of $-9,544$, but it simultaneously destroys general knowledge, evidenced by a retain rank collapse to $83,179$. In contrast, AdaPop provides a surgical intervention, achieving a strong unlearning signal (Rank $22,256$) while keeping the retain split stable (Rank $35.5$).

% \paragraph{AdaPop vs. WGA.} 
% Standard WGA fails to overcome the memory inertia of deeply embedded facts. On UNLamb (Llama-3), AdaPop achieves significantly deeper unlearning ($\Delta$ Log-Prob: $-33.4$) compared to WGA ($\Delta$ Log-Prob: $-9.2$). This confirms that fixed-weight ascent is insufficient for popular facts, whereas AdaPop’s adaptive exponent modulates the force required based on the fact's global footprint.

% \subsection{Model-Specific Observations}

% The architectural stability of the models varied under unlearning pressure:
% \begin{itemize}
%     \item \textbf{Llama-3.1-8B-Instruct:} The most robust architecture, exhibiting the cleanest separation between erasure and preservation.
%     \item \textbf{Gemma-7B-it:} Highly sensitive to unlearning; AdaPop effectively "confuses" the model on target facts without inducing linguistic collapse.
%     \item \textbf{Qwen-2.5-7B-Instruct:} Exhibited a "de-cluttering" effect. AdaPop not only achieved deep erasure ($-54$ $\Delta$ Log-Prob) but actually improved performance on the retain set. This suggests the dual-ascent controller acts as a corrective mechanism, sharpening the model's representations of the remaining knowledge.
% \end{itemize}

% \subsection{Results Discussion}

% Our results demonstrate that AdaPop is the superior method for real-world unlearning tasks. It achieves forgetting depths orders of magnitude better than standard WGA without the catastrophic risks associated with GA. The Retain Improvement Phenomenon suggests that by forcing the model to minimize retain loss while unlearning popular facts, the dual-ascent controller effectively "refreshes" the model's weights, leading to better-defined representations of the preserved knowledge.

\section{Conclusion}
\label{sec:conclusion}
TODO CONCLUSION
%% summary of result and that popularity is really important to check.



\section*{Limitations}
\textbf{Dependency on an external popularity proxy.}
AdaPop requires per-sample popularity scores. Wikidata \texttt{pop\_sum} provides good coverage for entity-centric factual QA but is unavailable for procedural or creative knowledge. LLM-estimated knowledgeability scores are a practical alternative, with correlation above $0.8$ with citation-based proxies~\cite{anna2026anatomy, marina2025llm}; the quality of the approximation directly limits the benefit of the popularity-dependent exponent.

\textbf{Calibration and evaluation scope.}
The exponent coefficients $a$ and $b$ were derived from the DUET distribution; applying AdaPop to proxies with a very different scale (e.g., PageRank) requires recalibration via the procedure in Appendix~\ref{app:derivation}. 
% All experiments use LoRA fine-tuning on 7--8B models; behaviour at larger scales, with full fine-tuning, or on non-entity knowledge types is not characterised.


\section*{Ethics Statement}
Our data come from public sources (HuggingFace, Wikidata, Wikipedia). We do not collect sensitive attributes, and no human subjects were involved. We view machine unlearning as a contribution to AI safety and data governance because it enables the removal of unsafe, outdated, or private content from deployed models. We also used ChatGPT 5 for minor language and grammatical edits; all research design, analysis, and interpretation were conducted by the authors.


% \section*{Acknowledgments}


% Bibliography entries for the entire Anthology, followed by custom entries
%\bibliography{anthology,custom}
% Custom bibliography entries only
\bibliography{custom}

\appendix

\label{sec:appendix}
\section{Derivation of AdaPop Popularity Parameters}

The power-law relationship $\beta_i = a \cdot s_i^{-b}$ is designed to map the wide dynamic range of popularity scores to a stable interval suitable for gradient updates. Since the gap between rare and popular entities typically spans 2--3 orders of magnitude in most popularity proxy metrics, the coefficients $a$ and $b$ must be tuned to ensure that unlearning remains effective yet controlled.

Specifically, we aim for "rare" entities (low $s_i$) to scale toward $\beta \approx 2.0$ to sharpen the update, while "popular" entities (high $s_i$) should scale toward $\beta \approx 0.05$ to flatten the update and overcome memory inertia.

\paragraph{Wikidata Proxy.} When the popularity proxy labeling is based on Wikidata (e.g., number of site links or claims), we empirically find the following coefficients to be optimal:
\begin{itemize}
    \item $a = 58.7$
    \item $b = 0.796$
\end{itemize}

These parameters may vary slightly depending on the specific range and distribution of the popularity proxy metric used (e.g., \texttt{pop\_sum} vs. \texttt{pagerank}). 

\paragraph{Parameter Derivation.} 
[Placeholder: Formal derivation of these coefficients based on the desired $\beta$ range and the cumulative distribution function (CDF) of the popularity scores.]

\section{Detailed UNLamb Results}

%% TALES OF 1E-3
% --- Table 1: Llama-3.1-8B-Instruct ---
\begin{table*}[h]
\centering
\small
\begin{tabularx}{\textwidth}{|l|Y|Y|Y|Y|Y|Y|Y|Y|}
\hline
\multirow{2}{*}{\textbf{Algorithm}} & \multicolumn{4}{c|}{\textbf{Popular Facts}} & \multicolumn{4}{c|}{\textbf{Rare Facts}} \\ \cline{2-9}
& \textbf{R.F.} $\downarrow$ & \textbf{R.R.} $\uparrow$ & \textbf{C.F.} $\downarrow$ & \textbf{C.R.} $\uparrow$ & \textbf{R.F.} $\downarrow$ & \textbf{R.R.} $\uparrow$ & \textbf{C.F.} $\downarrow$ & \textbf{C.R.} $\uparrow$ \\ \hline
\rowcolor{cyan!15} \textbf{Orig} & 0.820 & 0.827 & 0.491 & 0.555 & 0.821 & 0.827 & 0.518 & 0.555 \\ \hline
AdaPop & 0.274 & 0.889 & 0.047 & 0.866 & 0.253 & 0.902 & 0.189 & 0.868 \\ \hline
\rowcolor{gray!10} GA & 0.000 & 0.000 & 0.000 & 0.000 & 0.000 & 0.000 & 0.000 & 0.000 \\ \hline
GD & 0.000 & 0.000 & 0.000 & 0.000 & 0.000 & 0.018 & 0.000 & 0.012 \\ \hline
\rowcolor{gray!10} NPO & 0.363 & 0.881 & 0.092 & 0.873 & 0.346 & 0.876 & 0.174 & 0.809 \\ \hline
WGA & 0.239 & 0.902 & 0.159 & 0.900 & 0.275 & 0.914 & 0.200 & 0.871 \\ \hline
\end{tabularx}
\caption{Detailed UNLamb unlearning performance: Llama-3.1-8B-Instruct.}
\end{table*}

% --- Table 2: Qwen2.5-7B-Instruct ---
\begin{table*}[h]
\centering
\small
\begin{tabularx}{\textwidth}{|l|Y|Y|Y|Y|Y|Y|Y|Y|}
\hline
\multirow{2}{*}{\textbf{Algorithm}} & \multicolumn{4}{c|}{\textbf{Popular Facts}} & \multicolumn{4}{c|}{\textbf{Rare Facts}} \\ \cline{2-9}
& \textbf{R.F.} $\downarrow$ & \textbf{R.R.} $\uparrow$ & \textbf{C.F.} $\downarrow$ & \textbf{C.R.} $\uparrow$ & \textbf{R.F.} $\downarrow$ & \textbf{R.R.} $\uparrow$ & \textbf{C.F.} $\downarrow$ & \textbf{C.R.} $\uparrow$ \\ \hline
\rowcolor{cyan!15} \textbf{Orig} & 0.591 & 0.153 & 0.522 & 0.127 & 0.078 & 0.153 & 0.047 & 0.127 \\ \hline
AdaPop & 0.176 & 0.379 & 0.022 & 0.359 & 0.041 & 0.418 & 0.039 & 0.415 \\ \hline
\rowcolor{gray!10} GA & 0.000 & 0.000 & 0.000 & 0.000 & 0.000 & 0.000 & 0.000 & 0.000 \\ \hline
GD & 0.000 & 0.000 & 0.000 & 0.000 & 0.000 & 0.007 & 0.000 & 0.006 \\ \hline
\rowcolor{gray!10} NPO & 0.233 & 0.314 & 0.203 & 0.306 & 0.012 & 0.313 & 0.011 & 0.308 \\ \hline
WGA & 0.460 & 0.389 & 0.435 & 0.383 & 0.046 & 0.524 & 0.044 & 0.522 \\ \hline
\end{tabularx}
\caption{Detailed UNLamb unlearning performance: Qwen2.5-7B-Instruct.}
\end{table*}

% --- Table 3: Gemma-7B-it ---
\begin{table*}[h]
\centering
\small
\begin{tabularx}{\textwidth}{|l|Y|Y|Y|Y|Y|Y|Y|Y|}
\hline
\multirow{2}{*}{\textbf{Algorithm}} & \multicolumn{4}{c|}{\textbf{Popular Facts}} & \multicolumn{4}{c|}{\textbf{Rare Facts}} \\ \cline{2-9}
& \textbf{R.F.} $\downarrow$ & \textbf{R.R.} $\uparrow$ & \textbf{C.F.} $\downarrow$ & \textbf{C.R.} $\uparrow$ & \textbf{R.F.} $\downarrow$ & \textbf{R.R.} $\uparrow$ & \textbf{C.F.} $\downarrow$ & \textbf{C.R.} $\uparrow$ \\ \hline
\rowcolor{cyan!15} \textbf{Orig} & 0.683 & 0.283 & 0.087 & 0.089 & 0.156 & 0.283 & 0.044 & 0.089 \\ \hline
AdaPop & 0.244 & 0.422 & 0.057 & 0.383 & 0.040 & 0.459 & 0.030 & 0.444 \\ \hline
\rowcolor{gray!10} GA & 0.000 & 0.000 & 0.000 & 0.000 & 0.000 & 0.000 & 0.000 & 0.000 \\ \hline
GD & 0.000 & 0.000 & 0.000 & 0.000 & 0.000 & 0.006 & 0.000 & 0.000 \\ \hline
\rowcolor{gray!10} NPO & 0.110 & 0.201 & 0.014 & 0.149 & 0.010 & 0.160 & 0.005 & 0.138 \\ \hline
WGA & 0.288 & 0.429 & 0.110 & 0.397 & 0.030 & 0.454 & 0.024 & 0.450 \\ \hline
\end{tabularx}
\caption{Detailed UNLamb unlearning performance: Gemma-7B-it.}
\end{table*}


\end{document}
